% !TEX program = lualatex
% ---------------------------------------------------------------------------------------------
%  Readback note 37: The top floor at a finite place: finite Bessel functions, shell bounds and the top kernel
%
%  Build:  python3 tools/rb.py note <slug>
%
%  One unit of a Lean artifact, rewritten as ordinary mathematics. The prose is hand-written. Lean names,
%  line ranges and counts come from data.tex; quoted Lean statements are copied from the artifact by the build.
%
%  Rules of the text: standard terminology; every symbol is explained where it first appears and listed again
%  in the notation at the end; no letter has two meanings; a proof is complete or says what it takes for granted;
%  Lean names appear only in the LEAN lines and in the concordance.
%
%  The signature rule (the build enforces it): a function, relation or predicate that the artifact defines is
%  written only through its notation macro, which the build makes from the spec (`notation`) with one slot for
%  each explicit argument of the Lean definition. No argument is fixed and dropped, also not "for brevity" and
%  not for one section. Every argument and every hypothesis of a stated declaration stands inside its
%  statement, or in a Setting that the statement cites by number; only implicit and instance arguments (the
%  ambient objects) may stand in the conventions.
% ---------------------------------------------------------------------------------------------
\documentclass[11pt,a4paper]{article}
\ifdefined\pdfvariable\pdfvariable suppressoptionalinfo \numexpr 1+2+4+8\relax\fi
\usepackage[a4paper,hmargin=35mm,top=22mm,bottom=27mm,footskip=12mm,nohead]{geometry}

\usepackage{fontspec}
\directlua{
  luaotfload.add_fallback("gfallback", {"STIXTwoMath-Regular.otf:mode=node;", "DejaVuSans.ttf:mode=node;"})
  luaotfload.add_fallback("gmonofallback", {"DejaVuSansMono.ttf:mode=node;", "STIXTwoMath-Regular.otf:mode=node;"})
}
\setmainfont{NewCM10}[Extension=.otf,UprightFont=*-Book,ItalicFont=*-BookItalic,
  BoldFont=*-Bold,BoldItalicFont=*-BoldItalic,SmallCapsFont=*-Book,SmallCapsFeatures={Letters=SmallCaps},
  RawFeature={fallback=gfallback}]
\setsansfont{NewCMSans10}[Extension=.otf,UprightFont=*-Book,ItalicFont=*-BookOblique,
  BoldFont=*-Bold,BoldItalicFont=*-BoldOblique,RawFeature={fallback=gfallback}]
\setmonofont{JetBrainsMono}[Extension=.otf,UprightFont=*-Regular,ItalicFont=*-Italic,
  BoldFont=*-SemiBold,BoldItalicFont=*-SemiBoldItalic,Scale=MatchLowercase,
  RawFeature={-calt;-liga;fallback=gmonofallback}]
\usepackage{amsmath}
\usepackage{amsthm}
\usepackage[math-style=TeX]{unicode-math}
\setmathfont{NewCMMath-Book.otf}
\emergencystretch=2.5em\relax

\usepackage{xcolor}
\usepackage{tikz}
\usetikzlibrary{calc,arrows.meta,positioning,backgrounds}
\usepackage[protrusion=true,expansion=true]{microtype}
\usepackage{enumitem}
\usepackage{needspace}
\usepackage{booktabs}
\usepackage{array}
\usepackage{longtable}
\usepackage{fancyhdr}
\usepackage{lastpage}
\usepackage{multicol}
\usepackage[hidelinks,pdfusetitle]{hyperref}

\definecolor{ink}{HTML}{0F172A}
\definecolor{inksoft}{HTML}{475569}
\definecolor{inkfaint}{HTML}{94A3B8}
\definecolor{rulegrey}{HTML}{CBD5E1}
\definecolor{arith}{HTML}{C2410C}
\definecolor{est}{HTML}{7C3AED}
\definecolor{ana}{HTML}{0891B2}
\definecolor{ok}{HTML}{047857}
\definecolor{leancol}{HTML}{3730A3}

\title{The top floor at a finite place: finite Bessel functions, shell bounds and the top kernel}
\author{grh-compact}
\date{}

% ------------------------------------------------------------------------------------- macros
\newcommand{\leanfont}{\ttfamily\fontsize{8.6}{10.5}\selectfont\color{leancol}}
\newcommand{\lean}[1]{\csname LN@#1\endcsname}
\newcommand{\fillin}[1]{{\color{arith}\itshape[#1]}}  % a place of the template that is not written yet; the build reports every one that is left
\newcommand{\leanquote}[1]{\csname LQ@#1\endcsname}
% generated by rb.py note; do not edit
\newcommand{\ConcDecls}{45}
\newcommand{\Pin}{72fead10}
\newcommand{\PinText}{that of commit {\leanfont 72fead10}}
\newcommand{\ArtifactName}{grh-compact}
\expandafter\newcommand\csname LN@phase\endcsname{\href{https://github.com/lsernau/grh-lean/blob/72fead102f663d5c7a407964a13dc0cf9d977207/grh-compact/src/New/Filling/TopFloor.lean\#L23-L24}{{\leanfont IntegralPlace.\allowbreak{}phase}}}
\expandafter\newcommand\csname LN@existsbesselUofphasenezero\endcsname{\href{https://github.com/lsernau/grh-lean/blob/72fead102f663d5c7a407964a13dc0cf9d977207/grh-compact/src/New/Filling/TopFloor.lean\#L30-L31}{{\leanfont IntegralPlace.\allowbreak{}exists\_besselU\_of\_phase\_ne\_zero}}}
\expandafter\newcommand\csname LN@normphaseleone\endcsname{\href{https://github.com/lsernau/grh-lean/blob/72fead102f663d5c7a407964a13dc0cf9d977207/grh-compact/src/New/Filling/TopFloor.lean\#L33-L33}{{\leanfont IntegralPlace.\allowbreak{}norm\_phase\_le\_one}}}
\expandafter\newcommand\csname LN@phaseeq\endcsname{\href{https://github.com/lsernau/grh-lean/blob/72fead102f663d5c7a407964a13dc0cf9d977207/grh-compact/src/New/Filling/TopFloor.lean\#L35-L39}{{\leanfont IntegralPlace.\allowbreak{}phase\_eq}}}
\expandafter\newcommand\csname LN@phasebesselSmul\endcsname{\href{https://github.com/lsernau/grh-lean/blob/72fead102f663d5c7a407964a13dc0cf9d977207/grh-compact/src/New/Filling/TopFloor.lean\#L41-L50}{{\leanfont IntegralPlace.\allowbreak{}phase\_besselS\_mul}}}
\expandafter\newcommand\csname LN@phasemulofmemunitaryLevel\endcsname{\href{https://github.com/lsernau/grh-lean/blob/72fead102f663d5c7a407964a13dc0cf9d977207/grh-compact/src/New/Filling/TopFloor.lean\#L52-L53}{{\leanfont IntegralPlace.\allowbreak{}phase\_mul\_of\_mem\_unitaryLevel}}}
\expandafter\newcommand\csname LN@phasebesselUmul\endcsname{\href{https://github.com/lsernau/grh-lean/blob/72fead102f663d5c7a407964a13dc0cf9d977207/grh-compact/src/New/Filling/TopFloor.lean\#L55-L56}{{\leanfont IntegralPlace.\allowbreak{}phase\_besselU\_mul}}}
\expandafter\newcommand\csname LN@continuousphase\endcsname{\href{https://github.com/lsernau/grh-lean/blob/72fead102f663d5c7a407964a13dc0cf9d977207/grh-compact/src/New/Filling/TopFloor.lean\#L58-L60}{{\leanfont IntegralPlace.\allowbreak{}continuous\_phase}}}
\expandafter\newcommand\csname LN@unitaryLevelleunitaryLevelone\endcsname{\href{https://github.com/lsernau/grh-lean/blob/72fead102f663d5c7a407964a13dc0cf9d977207/grh-compact/src/New/Filling/TopFloor.lean\#L62-L63}{{\leanfont IntegralPlace.\allowbreak{}unitaryLevel\_le\_unitaryLevel\_one}}}
\expandafter\newcommand\csname LN@IntermediateGroup\endcsname{\href{https://github.com/lsernau/grh-lean/blob/72fead102f663d5c7a407964a13dc0cf9d977207/grh-compact/src/New/Filling/TopFloor.lean\#L69-L69}{{\leanfont Level.\allowbreak{}IntermediateGroup}}}
\expandafter\newcommand\csname LN@unipotentMass\endcsname{\href{https://github.com/lsernau/grh-lean/blob/72fead102f663d5c7a407964a13dc0cf9d977207/grh-compact/src/New/Filling/TopFloor.lean\#L73-L73}{{\leanfont Level.\allowbreak{}unipotentMass}}}
\expandafter\newcommand\csname LN@floorZeroForm\endcsname{\href{https://github.com/lsernau/grh-lean/blob/72fead102f663d5c7a407964a13dc0cf9d977207/grh-compact/src/New/Filling/TopFloor.lean\#L78-L79}{{\leanfont Level.\allowbreak{}floorZeroForm}}}
\expandafter\newcommand\csname LN@Cosets\endcsname{\href{https://github.com/lsernau/grh-lean/blob/72fead102f663d5c7a407964a13dc0cf9d977207/grh-compact/src/New/Filling/TopFloor.lean\#L82-L83}{{\leanfont Level.\allowbreak{}IntermediateGroup.\allowbreak{}Cosets}}}
\expandafter\newcommand\csname LN@massRatio\endcsname{\href{https://github.com/lsernau/grh-lean/blob/72fead102f663d5c7a407964a13dc0cf9d977207/grh-compact/src/New/Filling/TopFloor.lean\#L91-L92}{{\leanfont Level.\allowbreak{}IntermediateGroup.\allowbreak{}massRatio}}}
\expandafter\newcommand\csname LN@besselVector\endcsname{\href{https://github.com/lsernau/grh-lean/blob/72fead102f663d5c7a407964a13dc0cf9d977207/grh-compact/src/New/Filling/TopFloor.lean\#L96-L97}{{\leanfont Level.\allowbreak{}IntermediateGroup.\allowbreak{}besselVector}}}
\expandafter\newcommand\csname LN@unipotentMassnezeroandnetop\endcsname{\href{https://github.com/lsernau/grh-lean/blob/72fead102f663d5c7a407964a13dc0cf9d977207/grh-compact/src/New/Filling/TopFloor.lean\#L99-L100}{{\leanfont Level.\allowbreak{}unipotentMass\_ne\_zero\_and\_ne\_top}}}
\expandafter\newcommand\csname LN@existsofinnerbesselVectornezero\endcsname{\href{https://github.com/lsernau/grh-lean/blob/72fead102f663d5c7a407964a13dc0cf9d977207/grh-compact/src/New/Filling/TopFloor.lean\#L102-L109}{{\leanfont Level.\allowbreak{}exists\_of\_inner\_besselVector\_ne\_zero}}}
\expandafter\newcommand\csname LN@isOpen\endcsname{\href{https://github.com/lsernau/grh-lean/blob/72fead102f663d5c7a407964a13dc0cf9d977207/grh-compact/src/New/Filling/TopFloor.lean\#L111-L112}{{\leanfont Level.\allowbreak{}IntermediateGroup.\allowbreak{}isOpen}}}
\expandafter\newcommand\csname LN@existscoset\endcsname{\href{https://github.com/lsernau/grh-lean/blob/72fead102f663d5c7a407964a13dc0cf9d977207/grh-compact/src/New/Filling/TopFloor.lean\#L117-L119}{{\leanfont Level.\allowbreak{}IntermediateGroup.\allowbreak{}exists\_coset}}}
\expandafter\newcommand\csname LN@invmdeqsqrtgenIndexsq\endcsname{\href{https://github.com/lsernau/grh-lean/blob/72fead102f663d5c7a407964a13dc0cf9d977207/grh-compact/src/New/Filling/TopFloor.lean\#L121-L122}{{\leanfont IntegralPlace.\allowbreak{}inv\_md\_eq\_sqrt\_genIndex\_sq}}}
\expandafter\newcommand\csname LN@groupMeasurenetop\endcsname{\href{https://github.com/lsernau/grh-lean/blob/72fead102f663d5c7a407964a13dc0cf9d977207/grh-compact/src/New/Filling/TopFloor.lean\#L126-L127}{{\leanfont IntegralPlace.\allowbreak{}IsStdLevel.\allowbreak{}groupMeasure\_ne\_top}}}
\expandafter\newcommand\csname LN@massRationezeroandnetop\endcsname{\href{https://github.com/lsernau/grh-lean/blob/72fead102f663d5c7a407964a13dc0cf9d977207/grh-compact/src/New/Filling/TopFloor.lean\#L129-L131}{{\leanfont IntegralPlace.\allowbreak{}IsStdLevel.\allowbreak{}massRatio\_ne\_zero\_and\_ne\_top}}}
\expandafter\newcommand\csname LN@floorZeroFormoncongGroup\endcsname{\href{https://github.com/lsernau/grh-lean/blob/72fead102f663d5c7a407964a13dc0cf9d977207/grh-compact/src/New/Filling/TopFloor.lean\#L133-L138}{{\leanfont IntegralPlace.\allowbreak{}IsStdLevel.\allowbreak{}floorZeroForm\_on\_congGroup}}}
\expandafter\newcommand\csname LN@isCutbesselVector\endcsname{\href{https://github.com/lsernau/grh-lean/blob/72fead102f663d5c7a407964a13dc0cf9d977207/grh-compact/src/New/Filling/TopFloor.lean\#L140-L144}{{\leanfont IntegralPlace.\allowbreak{}IsStdLevel.\allowbreak{}isCut\_besselVector}}}
\expandafter\newcommand\csname LN@unipotentMassmulinnerbesselVector\endcsname{\href{https://github.com/lsernau/grh-lean/blob/72fead102f663d5c7a407964a13dc0cf9d977207/grh-compact/src/New/Filling/TopFloor.lean\#L148-L170}{{\leanfont IntegralPlace.\allowbreak{}IsStdLevel.\allowbreak{}unipotentMass\_mul\_inner\_besselVector}}}
\expandafter\newcommand\csname LN@pairingNormmullemassRatio\endcsname{\href{https://github.com/lsernau/grh-lean/blob/72fead102f663d5c7a407964a13dc0cf9d977207/grh-compact/src/New/Filling/TopFloor.lean\#L175-L197}{{\leanfont IntegralPlace.\allowbreak{}IsStdLevel.\allowbreak{}pairingNorm\_mul\_le\_massRatio}}}
\expandafter\newcommand\csname LN@enormintegralphaseequpperPairing\endcsname{\href{https://github.com/lsernau/grh-lean/blob/72fead102f663d5c7a407964a13dc0cf9d977207/grh-compact/src/New/Filling/TopFloor.lean\#L199-L205}{{\leanfont IntegralPlace.\allowbreak{}IsStdLevel.\allowbreak{}enorm\_integral\_phase\_eq\_upperPairing}}}
\expandafter\newcommand\csname LN@invtwopowlegroupMeasure\endcsname{\href{https://github.com/lsernau/grh-lean/blob/72fead102f663d5c7a407964a13dc0cf9d977207/grh-compact/src/New/Filling/TopFloor.lean\#L210-L213}{{\leanfont IntegralPlace.\allowbreak{}IsStdLevel.\allowbreak{}inv\_two\_pow\_le\_groupMeasure}}}
\expandafter\newcommand\csname LN@existsunitmemfloorValues\endcsname{\href{https://github.com/lsernau/grh-lean/blob/72fead102f663d5c7a407964a13dc0cf9d977207/grh-compact/src/New/Filling/TopFloor.lean\#L218-L234}{{\leanfont IntegralPlace.\allowbreak{}IsStdLevel.\allowbreak{}exists\_unit\_mem\_floorValues}}}
\expandafter\newcommand\csname LN@existscloseofinnernezero\endcsname{\href{https://github.com/lsernau/grh-lean/blob/72fead102f663d5c7a407964a13dc0cf9d977207/grh-compact/src/New/Filling/TopFloor.lean\#L238-L243}{{\leanfont IntegralPlace.\allowbreak{}IsStdLevel.\allowbreak{}exists\_close\_of\_inner\_ne\_zero}}}
\expandafter\newcommand\csname LN@existsgaugedepthofnezero\endcsname{\href{https://github.com/lsernau/grh-lean/blob/72fead102f663d5c7a407964a13dc0cf9d977207/grh-compact/src/New/Filling/TopFloor.lean\#L248-L269}{{\leanfont IntegralPlace.\allowbreak{}IsStdLevel.\allowbreak{}exists\_gauge\_depth\_of\_ne\_zero}}}
\expandafter\newcommand\csname LN@lintegralshellsle\endcsname{\href{https://github.com/lsernau/grh-lean/blob/72fead102f663d5c7a407964a13dc0cf9d977207/grh-compact/src/New/Filling/TopFloor.lean\#L275-L333}{{\leanfont IntegralPlace.\allowbreak{}IsStdLevel.\allowbreak{}lintegral\_shells\_le}}}
\expandafter\newcommand\csname LN@lintegralformMajorantnetop\endcsname{\href{https://github.com/lsernau/grh-lean/blob/72fead102f663d5c7a407964a13dc0cf9d977207/grh-compact/src/New/Filling/TopFloor.lean\#L336-L340}{{\leanfont IntegralPlace.\allowbreak{}IsStdLevel.\allowbreak{}lintegral\_formMajorant\_ne\_top}}}
\expandafter\newcommand\csname LN@enormfloorFormle\endcsname{\href{https://github.com/lsernau/grh-lean/blob/72fead102f663d5c7a407964a13dc0cf9d977207/grh-compact/src/New/Filling/TopFloor.lean\#L349-L382}{{\leanfont IntegralPlace.\allowbreak{}IsStdLevel.\allowbreak{}enorm\_floorForm\_le}}}
\expandafter\newcommand\csname LQ@enormfloorFormle\endcsname{\leanquoteline{0}{lemma enorm\_floorForm\_le \{ω\} (hK : 10 \textasciicircum{} 10 * 𝔰.n \textasciicircum{} 4 ≤ K)}\leanquoteline{4}{(H : (ω = fun x ↦ P.ψ (ξ * (T.frameGram x.1.1).diag ⟨j, by lia⟩)) ∨ ∃ i, j = i + 1 ∧ ω = 1 ∧ ξ =}\leanquoteline{4}{F.g⁻¹ * ((p i)⁻¹.1 - (q i)⁻¹.1)) : ‖F.floorForm j p q Z.besselVector ω‖ₑ ≤}\leanquoteline{2}{Z.massRatio * (.ofReal √(\_root\_.genIndex ξ.1) \textasciicircum{} t * (denomIndex R ξ.1 : ℝ≥0∞)⁻¹ \textasciicircum{} 6 *}\leanquoteline{2}{(2 / P.md F.g) \textasciicircum{} K * .ofReal (finHeadWeight K j p * finHeadWeight K j q))}}
\expandafter\newcommand\csname LN@floorZeroFormlemainterm\endcsname{\href{https://github.com/lsernau/grh-lean/blob/72fead102f663d5c7a407964a13dc0cf9d977207/grh-compact/src/New/Filling/TopFloor.lean\#L387-L415}{{\leanfont IntegralPlace.\allowbreak{}IsStdLevel.\allowbreak{}floorZeroForm\_le\_main\_term}}}
\expandafter\newcommand\csname LQ@floorZeroFormlemainterm\endcsname{\leanquoteline{0}{lemma floorZeroForm\_le\_main\_term : F.floorZeroForm Z ≤ F.groupMeasure Z * F.groupMeasure}\leanquoteline{2}{F.congGroup * (F.groupMeasure Z + volume (P.R \textbackslash{} P.intUnits) \textasciicircum{} (t + 1) * 2⁻¹ \textasciicircum{} 𝔰.n)}}
\expandafter\newcommand\csname LN@correlation\endcsname{\href{https://github.com/lsernau/grh-lean/blob/72fead102f663d5c7a407964a13dc0cf9d977207/grh-compact/src/New/Filling/TopFloor.lean\#L425-L426}{{\leanfont Level.\allowbreak{}correlation}}}
\expandafter\newcommand\csname LN@topKernel\endcsname{\href{https://github.com/lsernau/grh-lean/blob/72fead102f663d5c7a407964a13dc0cf9d977207/grh-compact/src/New/Filling/TopFloor.lean\#L430-L431}{{\leanfont Level.\allowbreak{}topKernel}}}
\expandafter\newcommand\csname LN@energya\endcsname{\href{https://github.com/lsernau/grh-lean/blob/72fead102f663d5c7a407964a13dc0cf9d977207/grh-compact/src/New/Filling/TopFloor.lean\#L436-L437}{{\leanfont Level.\allowbreak{}energy}}}
\expandafter\newcommand\csname LN@existsgaugeisHead\endcsname{\href{https://github.com/lsernau/grh-lean/blob/72fead102f663d5c7a407964a13dc0cf9d977207/grh-compact/src/New/Filling/TopFloor.lean\#L447-L451}{{\leanfont IntegralPlace.\allowbreak{}IsStdLevel.\allowbreak{}exists\_gauge\_isHead}}}
\expandafter\newcommand\csname LN@eventuallycorrelationeq\endcsname{\href{https://github.com/lsernau/grh-lean/blob/72fead102f663d5c7a407964a13dc0cf9d977207/grh-compact/src/New/Filling/TopFloor.lean\#L455-L469}{{\leanfont IntegralPlace.\allowbreak{}IsStdLevel.\allowbreak{}eventually\_correlation\_eq}}}
\expandafter\newcommand\csname LN@eventuallysumpacketsleenergy\endcsname{\href{https://github.com/lsernau/grh-lean/blob/72fead102f663d5c7a407964a13dc0cf9d977207/grh-compact/src/New/Filling/TopFloor.lean\#L474-L508}{{\leanfont IntegralPlace.\allowbreak{}IsStdLevel.\allowbreak{}eventually\_sum\_packets\_le\_energy}}}
\expandafter\newcommand\csname LQ@eventuallysumpacketsleenergy\endcsname{\leanquoteline{0}{lemma eventually\_sum\_packets\_le\_energy (hi : IsOpen \{x : A | IsUnit x\}) (hc :}\leanquoteline{2}{ContinuousAt Ring.inverse (1 : A)) \{i\} [Fintype i] \{l : i → \_\} (o : i) (hl : ∀ θ, l θ ∈ GT 𝔰.ϱ ∧}\leanquoteline{2}{orbitPoint 𝔰.top (l θ) = e) \{n : i → \_\} (hn : ∀ θ θ', l θ * (l θ')⁻¹ ∈ F.congGroup → n θ = n θ' →}\leanquoteline{2}{θ = θ') : let v θ p := correlation Z (ω p) (l θ) (n θ)}\leanquoteline{2}{∀ᶠ x in nhds e, ∀ k ∈ besselH ⊓ F.congGroup, ∀ y, k *ᵥ y = x →}\leanquoteline{12}{F.measureH F.congGroup * ∑ θ, .ofReal (star (v θ) ⬝ᵥ R *ᵥ v θ).re ≤ energy Z R ω y}}
\expandafter\newcommand\csname LN@floorFormtopeqintegral\endcsname{\href{https://github.com/lsernau/grh-lean/blob/72fead102f663d5c7a407964a13dc0cf9d977207/grh-compact/src/New/Filling/TopFloor.lean\#L513-L556}{{\leanfont IntegralPlace.\allowbreak{}IsStdLevel.\allowbreak{}floorForm\_top\_eq\_integral}}}
\expandafter\newcommand\csname LN@lintegralenergyletopKernel\endcsname{\href{https://github.com/lsernau/grh-lean/blob/72fead102f663d5c7a407964a13dc0cf9d977207/grh-compact/src/New/Filling/TopFloor.lean\#L561-L582}{{\leanfont IntegralPlace.\allowbreak{}IsStdLevel.\allowbreak{}lintegral\_energy\_le\_topKernel}}}
\expandafter\newcommand\csname LN@continuoustopKernel\endcsname{\href{https://github.com/lsernau/grh-lean/blob/72fead102f663d5c7a407964a13dc0cf9d977207/grh-compact/src/New/Filling/TopFloor.lean\#L586-L592}{{\leanfont IntegralPlace.\allowbreak{}IsStdLevel.\allowbreak{}continuous\_topKernel}}}
\expandafter\newcommand\csname LN@instFintypeCosets\endcsname{\href{https://github.com/lsernau/grh-lean/blob/72fead102f663d5c7a407964a13dc0cf9d977207/grh-compact/src/New/Filling/TopFloor.lean\#L85-L87}{{\leanfont instance : Fintype Z.\allowbreak{}Cosets}}}
\newcommand{\ConcordanceRows}{\crow{phase}{\href{https://github.com/lsernau/grh-lean/blob/72fead102f663d5c7a407964a13dc0cf9d977207/grh-compact/src/New/Filling/TopFloor.lean\#L23-L24}{{\leanfont IntegralPlace.\allowbreak{}phase}}}{TopFloor.lean}{23--24}
\crow{existsbesselUofphasenezero}{\href{https://github.com/lsernau/grh-lean/blob/72fead102f663d5c7a407964a13dc0cf9d977207/grh-compact/src/New/Filling/TopFloor.lean\#L30-L31}{{\leanfont IntegralPlace.\allowbreak{}exists\_besselU\_of\_phase\_ne\_zero}}}{TopFloor.lean}{30--31}
\crow{normphaseleone}{\href{https://github.com/lsernau/grh-lean/blob/72fead102f663d5c7a407964a13dc0cf9d977207/grh-compact/src/New/Filling/TopFloor.lean\#L33-L33}{{\leanfont IntegralPlace.\allowbreak{}norm\_phase\_le\_one}}}{TopFloor.lean}{33--33}
\crow{phaseeq}{\href{https://github.com/lsernau/grh-lean/blob/72fead102f663d5c7a407964a13dc0cf9d977207/grh-compact/src/New/Filling/TopFloor.lean\#L35-L39}{{\leanfont IntegralPlace.\allowbreak{}phase\_eq}}}{TopFloor.lean}{35--39}
\crow{phasebesselSmul}{\href{https://github.com/lsernau/grh-lean/blob/72fead102f663d5c7a407964a13dc0cf9d977207/grh-compact/src/New/Filling/TopFloor.lean\#L41-L50}{{\leanfont IntegralPlace.\allowbreak{}phase\_besselS\_mul}}}{TopFloor.lean}{41--50}
\crow{phasemulofmemunitaryLevel}{\href{https://github.com/lsernau/grh-lean/blob/72fead102f663d5c7a407964a13dc0cf9d977207/grh-compact/src/New/Filling/TopFloor.lean\#L52-L53}{{\leanfont IntegralPlace.\allowbreak{}phase\_mul\_of\_mem\_unitaryLevel}}}{TopFloor.lean}{52--53}
\crow{phasebesselUmul}{\href{https://github.com/lsernau/grh-lean/blob/72fead102f663d5c7a407964a13dc0cf9d977207/grh-compact/src/New/Filling/TopFloor.lean\#L55-L56}{{\leanfont IntegralPlace.\allowbreak{}phase\_besselU\_mul}}}{TopFloor.lean}{55--56}
\crow{continuousphase}{\href{https://github.com/lsernau/grh-lean/blob/72fead102f663d5c7a407964a13dc0cf9d977207/grh-compact/src/New/Filling/TopFloor.lean\#L58-L60}{{\leanfont IntegralPlace.\allowbreak{}continuous\_phase}}}{TopFloor.lean}{58--60}
\crow{unitaryLevelleunitaryLevelone}{\href{https://github.com/lsernau/grh-lean/blob/72fead102f663d5c7a407964a13dc0cf9d977207/grh-compact/src/New/Filling/TopFloor.lean\#L62-L63}{{\leanfont IntegralPlace.\allowbreak{}unitaryLevel\_le\_unitaryLevel\_one}}}{TopFloor.lean}{62--63}
\crow{IntermediateGroup}{\href{https://github.com/lsernau/grh-lean/blob/72fead102f663d5c7a407964a13dc0cf9d977207/grh-compact/src/New/Filling/TopFloor.lean\#L69-L69}{{\leanfont Level.\allowbreak{}IntermediateGroup}}}{TopFloor.lean}{69--69}
\crow{unipotentMass}{\href{https://github.com/lsernau/grh-lean/blob/72fead102f663d5c7a407964a13dc0cf9d977207/grh-compact/src/New/Filling/TopFloor.lean\#L73-L73}{{\leanfont Level.\allowbreak{}unipotentMass}}}{TopFloor.lean}{73--73}
\crow{floorZeroForm}{\href{https://github.com/lsernau/grh-lean/blob/72fead102f663d5c7a407964a13dc0cf9d977207/grh-compact/src/New/Filling/TopFloor.lean\#L78-L79}{{\leanfont Level.\allowbreak{}floorZeroForm}}}{TopFloor.lean}{78--79}
\crow{Cosets}{\href{https://github.com/lsernau/grh-lean/blob/72fead102f663d5c7a407964a13dc0cf9d977207/grh-compact/src/New/Filling/TopFloor.lean\#L82-L83}{{\leanfont Level.\allowbreak{}IntermediateGroup.\allowbreak{}Cosets}}}{TopFloor.lean}{82--83}
\crow{massRatio}{\href{https://github.com/lsernau/grh-lean/blob/72fead102f663d5c7a407964a13dc0cf9d977207/grh-compact/src/New/Filling/TopFloor.lean\#L91-L92}{{\leanfont Level.\allowbreak{}IntermediateGroup.\allowbreak{}massRatio}}}{TopFloor.lean}{91--92}
\crow{besselVector}{\href{https://github.com/lsernau/grh-lean/blob/72fead102f663d5c7a407964a13dc0cf9d977207/grh-compact/src/New/Filling/TopFloor.lean\#L96-L97}{{\leanfont Level.\allowbreak{}IntermediateGroup.\allowbreak{}besselVector}}}{TopFloor.lean}{96--97}
\crow{unipotentMassnezeroandnetop}{\href{https://github.com/lsernau/grh-lean/blob/72fead102f663d5c7a407964a13dc0cf9d977207/grh-compact/src/New/Filling/TopFloor.lean\#L99-L100}{{\leanfont Level.\allowbreak{}unipotentMass\_ne\_zero\_and\_ne\_top}}}{TopFloor.lean}{99--100}
\crow{existsofinnerbesselVectornezero}{\href{https://github.com/lsernau/grh-lean/blob/72fead102f663d5c7a407964a13dc0cf9d977207/grh-compact/src/New/Filling/TopFloor.lean\#L102-L109}{{\leanfont Level.\allowbreak{}exists\_of\_inner\_besselVector\_ne\_zero}}}{TopFloor.lean}{102--109}
\crow{isOpen}{\href{https://github.com/lsernau/grh-lean/blob/72fead102f663d5c7a407964a13dc0cf9d977207/grh-compact/src/New/Filling/TopFloor.lean\#L111-L112}{{\leanfont Level.\allowbreak{}IntermediateGroup.\allowbreak{}isOpen}}}{TopFloor.lean}{111--112}
\crow{existscoset}{\href{https://github.com/lsernau/grh-lean/blob/72fead102f663d5c7a407964a13dc0cf9d977207/grh-compact/src/New/Filling/TopFloor.lean\#L117-L119}{{\leanfont Level.\allowbreak{}IntermediateGroup.\allowbreak{}exists\_coset}}}{TopFloor.lean}{117--119}
\crow{invmdeqsqrtgenIndexsq}{\href{https://github.com/lsernau/grh-lean/blob/72fead102f663d5c7a407964a13dc0cf9d977207/grh-compact/src/New/Filling/TopFloor.lean\#L121-L122}{{\leanfont IntegralPlace.\allowbreak{}inv\_md\_eq\_sqrt\_genIndex\_sq}}}{TopFloor.lean}{121--122}
\crow{groupMeasurenetop}{\href{https://github.com/lsernau/grh-lean/blob/72fead102f663d5c7a407964a13dc0cf9d977207/grh-compact/src/New/Filling/TopFloor.lean\#L126-L127}{{\leanfont IntegralPlace.\allowbreak{}IsStdLevel.\allowbreak{}groupMeasure\_ne\_top}}}{TopFloor.lean}{126--127}
\crow{massRationezeroandnetop}{\href{https://github.com/lsernau/grh-lean/blob/72fead102f663d5c7a407964a13dc0cf9d977207/grh-compact/src/New/Filling/TopFloor.lean\#L129-L131}{{\leanfont IntegralPlace.\allowbreak{}IsStdLevel.\allowbreak{}massRatio\_ne\_zero\_and\_ne\_top}}}{TopFloor.lean}{129--131}
\crow{floorZeroFormoncongGroup}{\href{https://github.com/lsernau/grh-lean/blob/72fead102f663d5c7a407964a13dc0cf9d977207/grh-compact/src/New/Filling/TopFloor.lean\#L133-L138}{{\leanfont IntegralPlace.\allowbreak{}IsStdLevel.\allowbreak{}floorZeroForm\_on\_congGroup}}}{TopFloor.lean}{133--138}
\crow{isCutbesselVector}{\href{https://github.com/lsernau/grh-lean/blob/72fead102f663d5c7a407964a13dc0cf9d977207/grh-compact/src/New/Filling/TopFloor.lean\#L140-L144}{{\leanfont IntegralPlace.\allowbreak{}IsStdLevel.\allowbreak{}isCut\_besselVector}}}{TopFloor.lean}{140--144}
\crow{unipotentMassmulinnerbesselVector}{\href{https://github.com/lsernau/grh-lean/blob/72fead102f663d5c7a407964a13dc0cf9d977207/grh-compact/src/New/Filling/TopFloor.lean\#L148-L170}{{\leanfont IntegralPlace.\allowbreak{}IsStdLevel.\allowbreak{}unipotentMass\_mul\_inner\_besselVector}}}{TopFloor.lean}{148--170}
\crow{pairingNormmullemassRatio}{\href{https://github.com/lsernau/grh-lean/blob/72fead102f663d5c7a407964a13dc0cf9d977207/grh-compact/src/New/Filling/TopFloor.lean\#L175-L197}{{\leanfont IntegralPlace.\allowbreak{}IsStdLevel.\allowbreak{}pairingNorm\_mul\_le\_massRatio}}}{TopFloor.lean}{175--197}
\crow{enormintegralphaseequpperPairing}{\href{https://github.com/lsernau/grh-lean/blob/72fead102f663d5c7a407964a13dc0cf9d977207/grh-compact/src/New/Filling/TopFloor.lean\#L199-L205}{{\leanfont IntegralPlace.\allowbreak{}IsStdLevel.\allowbreak{}enorm\_integral\_phase\_eq\_upperPairing}}}{TopFloor.lean}{199--205}
\crow{invtwopowlegroupMeasure}{\href{https://github.com/lsernau/grh-lean/blob/72fead102f663d5c7a407964a13dc0cf9d977207/grh-compact/src/New/Filling/TopFloor.lean\#L210-L213}{{\leanfont IntegralPlace.\allowbreak{}IsStdLevel.\allowbreak{}inv\_two\_pow\_le\_groupMeasure}}}{TopFloor.lean}{210--213}
\crow{existsunitmemfloorValues}{\href{https://github.com/lsernau/grh-lean/blob/72fead102f663d5c7a407964a13dc0cf9d977207/grh-compact/src/New/Filling/TopFloor.lean\#L218-L234}{{\leanfont IntegralPlace.\allowbreak{}IsStdLevel.\allowbreak{}exists\_unit\_mem\_floorValues}}}{TopFloor.lean}{218--234}
\crow{existscloseofinnernezero}{\href{https://github.com/lsernau/grh-lean/blob/72fead102f663d5c7a407964a13dc0cf9d977207/grh-compact/src/New/Filling/TopFloor.lean\#L238-L243}{{\leanfont IntegralPlace.\allowbreak{}IsStdLevel.\allowbreak{}exists\_close\_of\_inner\_ne\_zero}}}{TopFloor.lean}{238--243}
\crow{existsgaugedepthofnezero}{\href{https://github.com/lsernau/grh-lean/blob/72fead102f663d5c7a407964a13dc0cf9d977207/grh-compact/src/New/Filling/TopFloor.lean\#L248-L269}{{\leanfont IntegralPlace.\allowbreak{}IsStdLevel.\allowbreak{}exists\_gauge\_depth\_of\_ne\_zero}}}{TopFloor.lean}{248--269}
\crow{lintegralshellsle}{\href{https://github.com/lsernau/grh-lean/blob/72fead102f663d5c7a407964a13dc0cf9d977207/grh-compact/src/New/Filling/TopFloor.lean\#L275-L333}{{\leanfont IntegralPlace.\allowbreak{}IsStdLevel.\allowbreak{}lintegral\_shells\_le}}}{TopFloor.lean}{275--333}
\crow{lintegralformMajorantnetop}{\href{https://github.com/lsernau/grh-lean/blob/72fead102f663d5c7a407964a13dc0cf9d977207/grh-compact/src/New/Filling/TopFloor.lean\#L336-L340}{{\leanfont IntegralPlace.\allowbreak{}IsStdLevel.\allowbreak{}lintegral\_formMajorant\_ne\_top}}}{TopFloor.lean}{336--340}
\crow{enormfloorFormle}{\href{https://github.com/lsernau/grh-lean/blob/72fead102f663d5c7a407964a13dc0cf9d977207/grh-compact/src/New/Filling/TopFloor.lean\#L349-L382}{{\leanfont IntegralPlace.\allowbreak{}IsStdLevel.\allowbreak{}enorm\_floorForm\_le}}}{TopFloor.lean}{349--382}
\crow{floorZeroFormlemainterm}{\href{https://github.com/lsernau/grh-lean/blob/72fead102f663d5c7a407964a13dc0cf9d977207/grh-compact/src/New/Filling/TopFloor.lean\#L387-L415}{{\leanfont IntegralPlace.\allowbreak{}IsStdLevel.\allowbreak{}floorZeroForm\_le\_main\_term}}}{TopFloor.lean}{387--415}
\crow{correlation}{\href{https://github.com/lsernau/grh-lean/blob/72fead102f663d5c7a407964a13dc0cf9d977207/grh-compact/src/New/Filling/TopFloor.lean\#L425-L426}{{\leanfont Level.\allowbreak{}correlation}}}{TopFloor.lean}{425--426}
\crow{topKernel}{\href{https://github.com/lsernau/grh-lean/blob/72fead102f663d5c7a407964a13dc0cf9d977207/grh-compact/src/New/Filling/TopFloor.lean\#L430-L431}{{\leanfont Level.\allowbreak{}topKernel}}}{TopFloor.lean}{430--431}
\crow{energya}{\href{https://github.com/lsernau/grh-lean/blob/72fead102f663d5c7a407964a13dc0cf9d977207/grh-compact/src/New/Filling/TopFloor.lean\#L436-L437}{{\leanfont Level.\allowbreak{}energy}}}{TopFloor.lean}{436--437}
\crow{existsgaugeisHead}{\href{https://github.com/lsernau/grh-lean/blob/72fead102f663d5c7a407964a13dc0cf9d977207/grh-compact/src/New/Filling/TopFloor.lean\#L447-L451}{{\leanfont IntegralPlace.\allowbreak{}IsStdLevel.\allowbreak{}exists\_gauge\_isHead}}}{TopFloor.lean}{447--451}
\crow{eventuallycorrelationeq}{\href{https://github.com/lsernau/grh-lean/blob/72fead102f663d5c7a407964a13dc0cf9d977207/grh-compact/src/New/Filling/TopFloor.lean\#L455-L469}{{\leanfont IntegralPlace.\allowbreak{}IsStdLevel.\allowbreak{}eventually\_correlation\_eq}}}{TopFloor.lean}{455--469}
\crow{eventuallysumpacketsleenergy}{\href{https://github.com/lsernau/grh-lean/blob/72fead102f663d5c7a407964a13dc0cf9d977207/grh-compact/src/New/Filling/TopFloor.lean\#L474-L508}{{\leanfont IntegralPlace.\allowbreak{}IsStdLevel.\allowbreak{}eventually\_sum\_packets\_le\_energy}}}{TopFloor.lean}{474--508}
\crow{floorFormtopeqintegral}{\href{https://github.com/lsernau/grh-lean/blob/72fead102f663d5c7a407964a13dc0cf9d977207/grh-compact/src/New/Filling/TopFloor.lean\#L513-L556}{{\leanfont IntegralPlace.\allowbreak{}IsStdLevel.\allowbreak{}floorForm\_top\_eq\_integral}}}{TopFloor.lean}{513--556}
\crow{lintegralenergyletopKernel}{\href{https://github.com/lsernau/grh-lean/blob/72fead102f663d5c7a407964a13dc0cf9d977207/grh-compact/src/New/Filling/TopFloor.lean\#L561-L582}{{\leanfont IntegralPlace.\allowbreak{}IsStdLevel.\allowbreak{}lintegral\_energy\_le\_topKernel}}}{TopFloor.lean}{561--582}
\crow{continuoustopKernel}{\href{https://github.com/lsernau/grh-lean/blob/72fead102f663d5c7a407964a13dc0cf9d977207/grh-compact/src/New/Filling/TopFloor.lean\#L586-L592}{{\leanfont IntegralPlace.\allowbreak{}IsStdLevel.\allowbreak{}continuous\_topKernel}}}{TopFloor.lean}{586--592}
\crow{instFintypeCosets}{\href{https://github.com/lsernau/grh-lean/blob/72fead102f663d5c7a407964a13dc0cf9d977207/grh-compact/src/New/Filling/TopFloor.lean\#L85-L87}{{\leanfont instance : Fintype Z.\allowbreak{}Cosets}}}{TopFloor.lean}{85--87}}
% the notations of the defined functions: one slot for each explicit argument of the Lean declaration
\newcommand{\phs}[2]{\chi_{#1}(#2)}
\newcommand{\IntG}[1]{\operatorname{Int}(#1)}
\newcommand{\umass}[1]{\mathsf{v}(#1)}
\newcommand{\fzf}[2]{\mathcal{I}_{0}(#1;#2)}
\newcommand{\Cos}[1]{\operatorname{Cos}(#1)}
\newcommand{\mr}[1]{\operatorname{mr}(#1)}
\newcommand{\bvec}[2]{W_{#1}(#2)}
\newcommand{\corr}[3]{\mathcal{J}(#1;#2,#3)}
\newcommand{\topK}[4]{\mathcal{K}(#1;#2;#3,#4)}
\newcommand{\UnitMembers}{497}
\newcommand{\UnitInNote}{45}
\newcommand{\UnitProofs}{179}
\newcommand{\UnitSteps}{1371}
\newcommand{\UnitAllSteps}{1479}


\linespread{1.08}                       % 11 pt on 14.7 pt
\setlength{\parindent}{0pt}
\setlength{\parskip}{5.5pt plus 1.5pt minus 0.5pt}
\setlength{\abovedisplayskip}{7.5pt plus 2pt minus 3pt}
\setlength{\belowdisplayskip}{7.5pt plus 2pt minus 3pt}
\setlength{\abovedisplayshortskip}{3pt plus 2pt}
\setlength{\belowdisplayshortskip}{6pt plus 2pt minus 2pt}
\emergencystretch=1.5em
\clubpenalty=10000 \widowpenalty=10000 \displaywidowpenalty=10000
\relpenalty=7000 \binoppenalty=9000
\raggedbottom
\pagestyle{fancy}\fancyhf{}\renewcommand{\headrulewidth}{0pt}
\fancyfoot[C]{\sffamily\fontsize{7.6}{9}\selectfont\color{inkfaint}\thepage\,/\,\pageref*{LastPage}}
\setlist{topsep=2pt,partopsep=0pt,parsep=0pt,itemsep=4pt,leftmargin=2.1em,labelsep=0.55em}

\newcommand{\Q}{\mathbb{Q}}
\newcommand{\Z}{\mathbb{Z}}
\newcommand{\F}{\mathbb{F}}
\newcommand{\len}{\operatorname{len}}
\newcommand{\cok}{\operatorname{cok}}
\newcommand{\End}{\operatorname{End}}
\newcommand{\Hom}{\operatorname{Hom}}
\newcommand{\Gal}{\operatorname{Gal}}
\newcommand{\diag}{\operatorname{diag}}
\newcommand{\kw}[1]{\emph{#1}}
\newcommand{\m}[1]{\mbox{$#1$}}        % a short formula that must not be broken across lines
\newcommand{\head}[1]{\par\addvspace{17pt plus 4pt minus 2pt}\Needspace{6\baselineskip}%
  \noindent{\fontsize{13.2}{16}\selectfont\bfseries\color{ink}#1}\par\nopagebreak\vspace{-1pt}}
\newcommand{\sect}[2]{\refstepcounter{section}\label{#2}\par\addvspace{19pt plus 4pt minus 2pt}\Needspace{7\baselineskip}%
  \noindent\llap{{\fontsize{13.2}{16}\selectfont\bfseries\color{inkfaint}\thesection}\hspace{9pt}}%
  {\fontsize{13.2}{16}\selectfont\bfseries\color{ink}#1}\par\nopagebreak\vspace{-1pt}}
\newcommand{\subhead}[1]{\par\addvspace{9pt plus 2pt}\Needspace{4\baselineskip}\noindent\textbf{#1}\enspace\ignorespaces}
% where a statement lives in the Lean archive
\newcommand{\leantag}{{\sffamily\bfseries\fontsize{6.3}{8}\selectfont\color{leancol}\addfontfeatures{LetterSpace=7}LEAN}}
\newcommand{\sep}{\unskip\nobreak\hspace{5pt}{\color{inkfaint}\textperiodcentered}\hspace{5pt plus 2pt}\ignorespaces}
\newcommand{\leanline}[1]{\par\nopagebreak\vspace{-2.5pt}{\fontsize{8.9}{12.4}\selectfont\color{inksoft}\raggedright
  \noindent\llap{\leantag\hspace{9pt}}#1\par}\vspace{1pt}}
% a displayed statement in a tinted box: #1 colour, #2 text
\newcommand{\statement}[2]{\par\addvspace{7pt}\begingroup\setlength{\fboxsep}{8.5pt}\setlength{\fboxrule}{0.6pt}%
  \noindent\fcolorbox{#1!30!white}{#1!3!white}{\parbox{\dimexpr\linewidth-2\fboxsep-2\fboxrule\relax}{#2}}%
  \endgroup\par\addvspace{7pt}}
\newcommand{\nt}[2]{\par\noindent\hangindent=1.1em\hangafter=1 #1\enspace{\color{inksoft}#2}\par\vspace{1pt}}
\newcommand{\legdot}[1]{\tikz[baseline=-0.62ex]\fill[#1] (0,0) circle (2.5pt);\hspace{3pt}}
\newcommand{\crow}[4]{\csname CS@#1\endcsname & #2 & \textcolor{inksoft}{#3} & \textcolor{inksoft}{#4}\\}
\newcommand{\leanquoteline}[2]{\par\noindent\hangindent=\dimexpr#1\fontdimen2\font+2.2em\relax\hangafter=1
  \hspace*{#1\fontdimen2\font}#2}
\newcommand{\quoted}[1]{\par\vspace{3pt}{\ttfamily\fontsize{7.5}{10.4}\selectfont\color{leancol}\raggedright
  \setlength{\parskip}{0pt}\leanquote{#1}\par}\vspace{4pt}}
\newcommand{\bx}[3]{{\sffamily\bfseries\fontsize{7.3}{9.2}\selectfont\color{#1!82!black}#2}\\[1.5pt]#3}
\newcommand{\cs}[2]{\expandafter\newcommand\csname CS@#1\endcsname{#2}}

% theorems
\newtheoremstyle{plainx}{7pt plus 2pt}{7pt plus 2pt}{\itshape}{}{\bfseries}{.}{0.5em}{}
\newtheoremstyle{uprightx}{7pt plus 2pt}{7pt plus 2pt}{}{}{\bfseries}{.}{0.5em}{}
\theoremstyle{plainx}
\newtheorem{theorem}{Theorem}[section]
\newtheorem{lemma}[theorem]{Lemma}
\newtheorem{proposition}[theorem]{Proposition}
\newtheorem{corollary}[theorem]{Corollary}
\theoremstyle{uprightx}
\newtheorem{definition}[theorem]{Definition}
\newtheorem{setting}[theorem]{Setting}          % hypotheses shared by several statements; each of them cites it by number
\newtheorem{remark}[theorem]{Remark}
\newtheorem{example}[theorem]{Example}
\renewcommand{\qedsymbol}{{\setlength{\fboxsep}{0pt}\setlength{\fboxrule}{0.5pt}\fbox{\rule{0pt}{0.46em}\hspace{0.46em}}}}
% what holds throughout the note: the ambient objects only (implicit and instance arguments), never a hypothesis
\newenvironment{conventions}{\par\addvspace{5pt}\noindent\textbf{Conventions.}\enspace\ignorespaces}{\par\addvspace{3pt}}
\makeatletter
\renewenvironment{proof}[1][Proof]{\par\pushQED{\qed}\normalfont\topsep0pt\partopsep0pt
  \trivlist\item[\hskip\labelsep\itshape #1.]\ignorespaces}{\popQED\endtrivlist\@endpefalse}
\makeatother

% where each declaration of the concordance is treated: \cs{key}{place}, the place naming a \label or a lettered item
\cs{phase}{Definition~\ref{d:phase}}
\cs{existsbesselUofphasenezero}{Lemma~\ref{l:phase}}
\cs{normphaseleone}{Lemma~\ref{l:phase}}
\cs{phaseeq}{Lemma~\ref{l:phase}}
\cs{phasebesselSmul}{Lemma~\ref{l:phase}}
\cs{phasemulofmemunitaryLevel}{Lemma~\ref{l:phase}}
\cs{phasebesselUmul}{Lemma~\ref{l:phase}}
\cs{continuousphase}{Lemma~\ref{l:phase}}
\cs{unitaryLevelleunitaryLevelone}{Lemma~\ref{l:phase}}
\cs{IntermediateGroup}{Definition~\ref{d:int}}
\cs{unipotentMass}{Definition~\ref{d:int}}
\cs{floorZeroForm}{Definition~\ref{d:int}}
\cs{Cosets}{Definition~\ref{d:int}}
\cs{massRatio}{Definition~\ref{d:int}}
\cs{besselVector}{Definition~\ref{d:int}}
\cs{unipotentMassnezeroandnetop}{Lemma~\ref{l:int}}
\cs{existsofinnerbesselVectornezero}{Lemma~\ref{l:inner}}
\cs{isOpen}{Lemma~\ref{l:int}}
\cs{existscoset}{Lemma~\ref{l:int}}
\cs{invmdeqsqrtgenIndexsq}{Lemma~\ref{l:int}}
\cs{groupMeasurenetop}{Lemma~\ref{l:int}}
\cs{massRationezeroandnetop}{Lemma~\ref{l:int}}
\cs{floorZeroFormoncongGroup}{Lemma~\ref{l:cong}}
\cs{isCutbesselVector}{Lemma~\ref{l:cut}}
\cs{unipotentMassmulinnerbesselVector}{Proposition~\ref{p:unfold}}
\cs{pairingNormmullemassRatio}{Proposition~\ref{p:pn}}
\cs{enormintegralphaseequpperPairing}{Lemma~\ref{l:avg}}
\cs{invtwopowlegroupMeasure}{Lemma~\ref{l:mass}}
\cs{existsunitmemfloorValues}{Lemma~\ref{l:val}}
\cs{existscloseofinnernezero}{Lemma~\ref{l:close}}
\cs{existsgaugedepthofnezero}{Proposition~\ref{p:depth}}
\cs{lintegralshellsle}{Proposition~\ref{p:shells}}
\cs{lintegralformMajorantnetop}{Corollary~\ref{c:finite}}
\cs{enormfloorFormle}{Theorem A}
\cs{floorZeroFormlemainterm}{Theorem B}
\cs{correlation}{Definition~\ref{d:top}}
\cs{topKernel}{Definition~\ref{d:top}}
\cs{energya}{Definition~\ref{d:top}}
\cs{existsgaugeisHead}{Lemma~\ref{l:heads}}
\cs{eventuallycorrelationeq}{Proposition~\ref{p:const}}
\cs{eventuallysumpacketsleenergy}{Theorem D}
\cs{floorFormtopeqintegral}{Theorem C}
\cs{lintegralenergyletopKernel}{Theorem D}
\cs{continuoustopKernel}{Theorem C}
\cs{instFintypeCosets}{Definition~\ref{d:int}}
% ----- macros-148: the recalled notions of unit 148 (chunks a, b) and of its neighbours, with every explicit
% ----- argument of the Lean definition as a slot. Shared by Notes F, P, T, G; copy unchanged into the preamble.
\newcommand{\fs}{\mathfrak{s}}                      % the shape (rho, kappa, k_0)
\newcommand{\sT}{\mathsf{T}}                        % a finite set of finite places
\newcommand{\cL}{\mathcal{L}}                       % a level
\newcommand{\AEf}{\mathbb{A}_{E,f}}
\newcommand{\hO}{\hat{O}_E}
\newcommand{\GL}{\mathrm{GL}}
\newcommand{\full}{\mathrm{full}}
\newcommand{\cut}{\mathrm{cut}}
\newcommand{\rh}[1]{\varrho(#1)}                    % Shape.ϱ
\newcommand{\kk}[1]{k(#1)}                          % Shape.k = #κ
\newcommand{\nn}[1]{n(#1)}                          % Shape.n = ϱ + 1
\newcommand{\rr}[1]{r(#1)}                          % Shape.r = 2ϱ + k + 1
\newcommand{\Ix}[1]{I(#1)}                          % Shape.BesselIndex
\newcommand{\Mx}[1]{M(#1)}                          % Shape.MiddleIndex
\newcommand{\ninf}[1]{f(#1)}                        % numInfPlaces E
\newcommand{\ET}[1]{E_{#1}}                         % LocProd T
% the components of the filling data S
\newcommand{\bS}[1]{b(#1)}                          % S.b
\newcommand{\NS}[1]{N(#1)}                          % S.N
\newcommand{\qS}[2]{q_{#2}(#1)}                     % S.levelPrime i
\newcommand{\gS}[1]{g(#1)}                          % S.g = q_0 q_1
\newcommand{\piS}[1]{\bar\pi(#1)}                   % S.piBar
\newcommand{\vpS}[1]{\varpi(#1)}                    % S.ϖ
\newcommand{\ChS}[2]{\mathrm{Ch}_{#2}(#1)}          % S.chosenPlaces i
\newcommand{\KcolS}[1]{K^{\mathrm{col}}(#1)}        % S.columnLevel
\newcommand{\BsysS}[2]{\mathfrak{B}_{#2}(#1)}       % S.besselSys v
% forms, groups, levels of the set-up (Setup.lean)
\newcommand{\Trat}[2]{T^{\mathrm{rat}}_{#1}(#2)}    % S.ratMiddle R
\newcommand{\Tlat}[2]{T^{\mathrm{lat}}_{#1}(#2)}    % S.latMiddle R
\newcommand{\GS}[2]{G_{#1}(#2)}                     % S.G R
\newcommand{\GFin}[1]{G^{\mathrm{lat}}_{#1}}        % S.GFin
\newcommand{\GRes}[1]{\mathsf{G}(#1)}               % S.GRes
\newcommand{\Cl}[1]{\mathrm{Cl}(#1)}                % S.Classes
\newcommand{\Bc}[1]{\mathrm{Bc}(#1)}                % S.BesselCoord
\newcommand{\KS}[1]{K(#1)}                          % S.latStabilizer
\newcommand{\redS}[1]{\mathrm{red}_{#1}}            % S.reduction
\newcommand{\KgS}[1]{K_g(#1)}                       % S.principalLevel
\newcommand{\redim}[2]{\overline{#2}^{\,#1}}        % S.reducedImage A
\newcommand{\kapS}[2]{\kappa_{#2}(#1)}              % S.kappa v
\newcommand{\gid}[1]{\hat{g}(#1)}                   % S.gIdele
\newcommand{\latB}[1]{\Lambda(#1)}                  % S.latBasis
\newcommand{\rinf}[1]{\mathrm{r}^{\infty}_{#1}}     % S.ratToInf
\newcommand{\rfin}[1]{\mathrm{r}^{f}_{#1}}          % S.ratToFin
\newcommand{\incH}[1]{\mathrm{inc}_{#1}}            % S.inclHFin
\newcommand{\hc}[2]{h_{#1}(#2)}                     % S.h c
\newcommand{\pos}[4]{\mathrm{pos}_{#1}(#2,#3;#4)}   % S.termPosition c c' γ
\newcommand{\terms}[4]{\mathrm{Tm}_{#1}(#2,#3,#4)}  % S.termsOver c s c'
\newcommand{\Zf}[2]{Z_{#1}(#2)}                     % S.frameProducts γ
\newcommand{\tdisc}[2]{\mathsf{g}_{#1}(#2)}         % S.termDisc γ
\newcommand{\clU}[2]{\mathcal{U}_{#1,#2}}           % S.classUnip c
\newcommand{\cutoff}[2]{\beta_{#1,#2}}              % S.classCutoff c
\newcommand{\besfn}[4]{W^{#1}_{#2}(#3)_{#4}}        % S.besselFn y i, with the side
\newcommand{\amp}[5]{a^{#1}_{#2}(#3,#4;#5)}         % S.amplitude c c' γ, with the side
\newcommand{\ampabs}[5]{|a|^{#1}_{#2}(#3,#4;#5)}    % S.amplitudeAbs c c' γ, with the side
\newcommand{\ampsum}[6]{A^{#1}_{#2}(#3;#4,#5,#6)}   % S.amplitudeSum T c s c', with the side
\newcommand{\uchar}[2]{\psi^{\mathcal{U}}_{#1}(#2)} % S.unipotentChar u
\newcommand{\ucov}[1]{\nu(#1)}                      % S.unipCovolume
\newcommand{\siegC}[1]{c_{*}(#1)}                   % S.siegelConst
\newcommand{\packC}[1]{a_{*}(#1)}                   % S.packetConst
\newcommand{\uord}[1]{\mathrm{uo}(#1)}              % S.unipOrder
\newcommand{\cdim}[1]{\mathrm{cd}(#1)}              % S.columnDim
\newcommand{\bdeg}[1]{\deg(#1)}                     % S.besselDegree
\newcommand{\mainE}[1]{\mathsf{d}(#1)}              % S.mainEnergy
\newcommand{\besE}[1]{\mathsf{d}_{1}(#1)}           % S.besselEnergy
\newcommand{\tailS}[1]{\mathrm{tl}(#1)}             % S.tailScale
\newcommand{\codimF}[1]{\mathrm{cf}(#1)}            % S.codimFactor
\newcommand{\fscale}[3]{\mathrm{fs}_{#1}(#2,#3)}    % S.frameScaleSq v a
\newcommand{\fsprod}[1]{\mathrm{fs}(#1)}            % S.frameScaleProd
\newcommand{\mass}[2]{\operatorname{mass}(#1,#2)}   % mass E r
% chunk b and the sieve
\newcommand{\energy}[3]{\mathrm{En}_{#1}(#2;#3)}    % Filling.energy A Ω (S implicit in Lean, shown)
\newcommand{\linkE}[5]{\mathrm{LE}^{#1}_{#2}(#3;#4,#5)} % Filling.linkEnergy T Ω' R (side and S implicit, shown)
\newcommand{\fixset}[3]{\mathrm{Fix}_{#1}(#2,#3)}   % Filling.fixingSet Sg d (S implicit, shown)
\newcommand{\vfix}[3]{\mathrm{Fix}_{#1,#2}(#3)}     % Filling.vectorFixer w d (S implicit, shown)
\newcommand{\charT}[2]{\psi_{#1,#2}}                % S.charOnPlaces T
\newcommand{\fgain}[3]{\beta_{#1,#2}(#3)}           % S.frequencyGain T ξ
\newcommand{\link}[1]{\operatorname{link}(#1)}      % link γ
\newcommand{\clink}[2]{\operatorname{link}^{\circ}(#1,#2)} % commonLink γ γ'
\newcommand{\tdnorm}[2]{\|\mathsf{g}\|_{#1}(#2)}    % S.termDiscNorm γ
\newcommand{\mnorm}[3]{\mathrm{mn}_{#1}(#2;#3)}     % S.minorNorm j γ
\newcommand{\cgauge}[5]{\mathrm{cg}_{#1}(#2;#3,#4;#5)} % S.classGauge j c c' γ
\newcommand{\pinv}[5]{\mathrm{pi}_{#1}(#2,#3;#4,#5)} % S.pairInvariant c c' γ γ'
\newcommand{\ldsq}[3]{\mathrm{ld}_{#1}(#2,#3)}      % S.linkDiscSq v e
\newcommand{\lmet}[1]{Q^{\mathrm{link}}_{#1}}       % S.linkMetric
% ----- end of macros-148
% ----- notions of the background library, of the unit on the real place and of the note on standard levels that
% ----- this note recalls, written by hand with every explicit argument (TeX as in the notes that state them)
\newcommand{\JT}[1]{J(#1)}                          % Shape.Form.splitGram
\newcommand{\GT}[1]{G(#1)}                          % Shape.Form.G
\newcommand{\UT}[1]{\mathcal{U}(#1)}                % besselU
\newcommand{\HT}[1]{\mathcal{H}(#1)}                % besselH
\newcommand{\ST}[1]{\mathcal{S}(#1)}                % besselS
\newcommand{\Hl}[2]{\mathcal{H}_{#1}(#2)}           % GT l
\newcommand{\Ql}[2]{\mathcal{Q}_{#1}(#2)}           % unitriGroup l
\newcommand{\sd}[1]{\operatorname{sd}(#1)}          % simpleRootSum
\newcommand{\ZT}[2]{Z_{#1}(#2)}                     % Shape.Form.frameGram
\newcommand{\weyl}[2]{w_{#1}(#2)}                   % weylElem l p
\newcommand{\gix}[1]{[#1]}                          % genIndex
\newcommand{\den}[1]{\operatorname{den}(#1)}        % denomIndex
\newcommand{\modA}{\operatorname{mod}_A}            % the modulus of the place structure
\newcommand{\zetaA}[1]{\zeta_{A}(#1)}               % idealZeta O_A
\newcommand{\uLev}[2]{K_{#2}(#1)}                   % unitaryLevel O_A g J(T)
\newcommand{\cgp}[1]{\mathrm{K}^{\circ}(#1)}        % Level.congGroup F
\newcommand{\gL}[1]{g_{#1}}                         % the unit of a level
\newcommand{\BoxL}[1]{\operatorname{Box}(#1)}       % boxSet
\newcommand{\Posn}[1]{\operatorname{Posn}(#1)}      % positions
\newcommand{\Levels}[3]{\operatorname{Lev}(#1,#2,#3)} % Level ι T m
\newcommand{\Gauges}[1]{\operatorname{Gau}(#1)}     % Gauge F
\newcommand{\IsStd}[1]{#1\in\operatorname{StdLev}}  % IntegralPlace.IsStdLevel F
\newcommand{\stdg}[2]{\Gamma_{#1}(#2)}              % IntegralPlace.IsStdLevel.gauge F α
\newcommand{\Lam}[1]{\Lambda(#1)}                   % Gauge.Λ
\newcommand{\szl}[3]{\operatorname{sz}_{#1}(#2)\le #3} % x ∈ Γ.ofSize i, Γ.HasMatSize i X
\newcommand{\Hd}[3]{#3\in\operatorname{Hd}_{#2}(#1)} % Γ.IsHead j p
\newcommand{\gM}[1]{\mathsf{M}(#1)}                 % Level.groupMeasure F
\newcommand{\gMu}[1]{\mathsf{M}_{\ge0}(#1)}         % Level.upperMeasure F 0
\newcommand{\flm}[2]{\theta_{#2}(#1)}               % Level.floorMeasure F l
\newcommand{\mH}[1]{\theta_{n}(#1)}                 % Level.measureH F
\newcommand{\stk}[3]{\mu^{#1}_{#2,#3}}              % Level.stackMeasure F l i
\newcommand{\quadric}[2]{\delta^{(#2)}(#1)}         % floorQuadric
\newcommand{\etab}[1]{\eta_{<#1}}                   % ascConv heisenbergMeasure j
\newcommand{\fmod}[3]{\mathsf{d}_{#2}(#1;#3)}       % Level.floorModulus F l g
\newcommand{\vexp}[1]{\mathsf{s}(#1)}               % Level.volumeExponent F
\newcommand{\Val}[4]{\operatorname{Val}_{#2}(#1;#3,#4)} % Level.floorValues F a x y
\newcommand{\Av}[5]{\mathsf{Av}_{#2}(#1;#3)(#4,#5)} % Level.upperPairing F j k g y
\newcommand{\Pv}[5]{\mathsf{Pv}_{#2}(#1;#3)(#4,#5)} % Level.pivotModulus F j k g y
\newcommand{\Sm}[5]{\mathsf{Sm}_{#2}(#1;#3,#4)(#5)} % Level.weightedSquareSum F k x j g
\newcommand{\PN}[3]{\operatorname{PN}(#1;#2,#3)}    % Level.pairingNorm F W W'
\newcommand{\hdens}[4]{\mathsf{h}_{#2}(#1;#3,#4)}   % Level.headDensity F j p q
\newcommand{\fform}[6]{\mathcal{I}_{#2}(#1;#3,#4;#5,#6)} % Level.floorForm F j p q W ω
\newcommand{\fmodu}[7]{\mathcal{I}^{\mathrm{mod}}_{#2}(#1;#3,#4;#5,#6,#7)} % Level.formModulus F j p q W x ω
\newcommand{\fmaj}[7]{\mathcal{I}^{\mathrm{maj}}_{#2}(#1;#3,#4;#5,#6,#7)} % Level.formMajorant F j p q W W' x
\newcommand{\snp}[7]{\operatorname{SNP}_{#2}(#1;#3,#4;#5,#6,#7)} % Gauge.selfNormProduct Γ j p q W W' x
\newcommand{\bcorr}[3]{\mathrm{J}(#1,#2;#3)}        % besselCorrelation μ W γ
\newcommand{\fhw}[3]{\Lambda_{#1}(#2;#3)}           % finHeadWeight K i p
\newcommand{\lenergy}[4]{\mathcal{E}(#1;#2,#3,#4)}  % Level.energy Z R ω e (Notes on the formal text, 1)
\newcommand{\pk}[1]{\mathsf{k}_{#1}}                % the kernel (x,y) -> <W_Z(x), W_Z(y)> of an intermediate group
\newcommand{\vol}{\operatorname{vol}}
\newcommand{\rePart}{\operatorname{Re}}
\newcommand{\CC}{\mathbb{C}}
\newcommand{\RR}{\mathbb{R}}
\newcommand{\NN}{\mathbb{N}}
\newcommand{\Inp}[1]{Input~\ref{i:#1}}

\begin{document}
\color{ink}

{\raggedright
{\sffamily\bfseries\fontsize{7.2}{9}\selectfont\color{leancol}\addfontfeatures{LetterSpace=8}READBACK NOTE 37\par}
\vspace{5pt}
{\fontsize{21.5}{25.5}\selectfont\bfseries The top floor at a finite place: finite Bessel functions, shell bounds and the top kernel\par}
\vspace{5pt}
{\fontsize{12.4}{16}\selectfont\itshape\color{inksoft} One unit of the Lean artifact \ArtifactName, read back into
mathematics\par}}
\vspace{8pt}
{\color{rulegrey}\hrule height 0.6pt}
\vspace{7pt}
\thispagestyle{empty}

\textbf{Summary.} At finitely many finite places the second moment of the Poincaré series is governed by one
function on a unitary group over a product $A$ of local fields: the \emph{finite Bessel function}, which is a
character of the Bessel unipotent subgroup on the cosets of a principal congruence subgroup and zero elsewhere.
This note feeds the vector of its translates, the \emph{finite Bessel vector} of a group $Z$ between the
congruence subgroup and the integral group, into the tower method of the unit on the real place. It proves three
things. The floor forms of this vector decay in the frequency: a non-integral frequency~$\xi$ costs a power of
the index of~$\xi$. The form at the bottom of the tower is the square of the measure of~$Z$ up to an error that is
a high power of the volume of the non-units. And at the top floor the form unfolds to an integral of inner
products of \emph{correlations}, which gives a kernel that is continuous, of positive type, and bounds from
above an energy that counts packets of rational terms. The estimates rest on one \emph{shell bound}: the
positions are sorted by the denominator ideal of $(x-1)/g$, and a sum over ideals is paid for by an ideal zeta
function.

\statement{leancol}{\fontsize{9.6}{13.2}\selectfont\textbf{What this note is.} The artifact \ArtifactName{} is a
body of Lean code: the compact edition of a formal proof, due to L.~Alpöge with the AI system Claude, of the
Riemann hypothesis and of its generalisation to Dirichlet $L$-functions. That headline is a claim. The authors
report that every file compiles and that Lean's kernel accepted the proofs; an independent replay is recorded in
the handout's record of checks when it reports. No mathematician has read the argument through, and nothing here
says that the Riemann hypothesis is proved. This note (Note~T of its unit) treats \UnitInNote{} of the
\UnitMembers{} declarations of the unit \emph{Translates of Poincaré series fill the rows}: the module on the
top floor at a finite place. Proved here from the inputs: Theorems A to~D and the statements of
§§\ref{s:phase}--\ref{s:top}. Taken for granted and stated as Inputs 1 to~5 in §\ref{s:inputs}: the Bessel
subgroup and the congruence subgroups from the background library, the tower of floor measures and its
estimates from the unit on the real place, the standard levels from another note of this unit, and three facts
on ideals, valuations and Eichler elements. Every statement was compared with its Lean text, argument by
argument. The Lean text read is \PinText; Appendix~A is the concordance of \ConcDecls{} declarations.}

\sect{The statements}{s:intro}
The objects are defined in §\ref{s:setting}. In brief: $A$ is a finite product of local fields with an
involution, $T$ a hermitian form of a shape~$\fs$, \m{\GT{T}} the unitary group of its split Gram matrix, $\cL$ a
standard level with unit \m{\gL{\cL}} and congruence subgroup \m{\cgp{\cL}}, and $Z$ a group between
\m{\cgp{\cL}} and the group \m{\uLev{1}{T}} of the integral isometries. The finite Bessel vector \m{\bvec{Z}{y}}
has one coordinate for each coset of \m{\cgp{\cL}} in~$Z$; \m{\fform{\cL}{j}{p}{q}{W}{\omega}} is the floor form of
the unit on the real place, \m{\mr{Z}} the mass ratio, \m{\gix{\xi}} the generalised index and
\m{\den{\xi}} the denominator index of a unit~$\xi$.

\statement{arith}{\phantomsection\label{t:A}\textbf{Theorem A} (the frequency bound).\enspace
\emph{In Setting~\ref{s:Z}, let \m{j\le\rh{\fs}}, let $K$ and $t$ be natural numbers with
\m{10^{10}\nn{\fs}^4\le K} and \m{t+10^{10}\nn{\fs}^4\le\kk{\fs}}, let $\xi$ be a unit of~$A$, let $p$, $q$ be
sequences of units of~$A$, and let \m{\omega\colon\GT{T}\to\CC} be a function such that either
\m{\omega(x)=\psi_A\bigl(\xi\cdot\ZT{T}{x}_{jj}\bigr)} for all~$x$, or there is $i$ with \m{j=i+1}, \m{\omega=1} and
\m{\xi=\gL{\cL}^{-1}(p_i^{-1}-q_i^{-1})}. Then}
\[\bigl|\fform{\cL}{j}{p}{q}{\bvec{Z}{\cdot}}{\omega}\bigr|\ \le\ \mr{Z}\cdot\gix{\xi}^{t/2}\,\den{\xi}^{-6}\,
\Bigl(\frac{2}{\modA(\gL{\cL})}\Bigr)^{K}\,\fhw{K}{j}{p}\,\fhw{K}{j}{q}.\]}
\leanline{\lean{enormfloorFormle}}

\statement{arith}{\phantomsection\label{t:B}\textbf{Theorem B} (the bottom value).\enspace
\emph{In Setting~\ref{s:Z}, let $t$ be a natural number with \m{t+10^{10}\nn{\fs}^4\le\kk{\fs}}. Then}
\[\fzf{\cL}{Z}\ \le\ \gM{\cL}(Z)\cdot\gM{\cL}\bigl(\cgp{\cL}\bigr)\cdot\Bigl(\gM{\cL}(Z)+\vol\bigl(O_A\setminus
O_A^{\times}\bigr)^{t+1}\cdot2^{-\nn{\fs}}\Bigr).\]}
\leanline{\lean{floorZeroFormlemainterm}}

\statement{arith}{\phantomsection\label{t:C}\textbf{Theorem C} (the top floor: unfolding and continuity).\enspace
\emph{In Setting~\ref{s:Z}:}
\begin{enumerate}[label=\textup{(\roman*)},leftmargin=2.3em,itemsep=1pt]
\item \emph{Let $p$, $q$ be sequences of units of~$A$ and \m{\xi\in A}; put
\m{\varphi_\xi(x)=\psi_A\bigl(\xi\cdot\ZT{T}{x}_{\varrho\varrho}\bigr)} for \m{x\in\GT{T}}. Then the function
\m{F_\xi(l)=\varphi_\xi(l^{-1})\cdot\langle\corr{Z}{p}{l},\corr{Z}{q}{l}\rangle} is integrable for the measure
\m{\stk{\cL}{\varrho}{2}}, and}
\[\fform{\cL}{\varrho}{p}{q}{\bvec{Z}{\cdot}}{\varphi_\xi}\ =\ \int F_\xi(l)\,d\stk{\cL}{\varrho}{2}(l).\]
\item \emph{Let $L$ be an open additive subgroup of~$A$. Then \m{(p,q)\mapsto\topK{L}{Z}{p}{q}} is a continuous
function on the pairs of sequences of units of~$A$.}
\end{enumerate}}
\leanline{\lean{floorFormtopeqintegral}\sep \lean{continuoustopKernel}}

\statement{arith}{\phantomsection\label{t:D}\textbf{Theorem D} (the top floor: energy and packets).\enspace
\emph{In Setting~\ref{s:Z}, let $\mathcal C$ be a finite set, $R$ a complex matrix with rows and columns indexed
by~$\mathcal C$, and \m{\omega=(\omega_i)_{i\in\mathcal C}} a family of sequences of units of~$A$.}
\begin{enumerate}[label=\textup{(\roman*)},leftmargin=2.3em,itemsep=1pt]
\item \emph{Let \m{L\subseteq A} be an open subset, suppose that $R$ is positive semidefinite, and let $t$ be a
unit of $A$ with \m{(\omega_i)_{\varrho}=t} for all indices \m{i\in\mathcal C}. Then
\m{\mathcal Q=\sum_{i,i'\in\mathcal C}R_{ii'}\cdot\topK{L}{Z}{\omega_i}{\omega_{i'}}} is real and \m{\mathcal Q\ge0}, and}
\[\int\mathbf 1\bigl[(\JT{T}\,e)_{z}-t\in L\bigr]\cdot\lenergy{Z}{R}{\omega}{e}\ d\quadric{\cL}{\varrho}(e)\ \le\ \mathcal Q.\]
\item \emph{Suppose that the units of $A$ form an open set and that inversion is continuous at~$1$. Let $e$ be a
vector, $\Theta_0$ a non-empty finite set, and for \m{\theta\in\Theta_0} let \m{l_\theta\in\Hl{\varrho}{T}} with
\m{l_\theta^{-1}z=e} and \m{n_\theta\in\Cos{Z}\times\Cos{Z}}, such that \m{l_\theta l_{\theta'}^{-1}\in\cgp{\cL}} and
\m{n_\theta=n_{\theta'}} together imply \m{\theta=\theta'}. For \m{\theta\in\Theta_0} let $v_\theta$ be the vector with the
coordinate \m{\corr{Z}{\omega_i}{l_\theta}_{n_\theta}} for every index \m{i\in\mathcal C}. Then there is a neighbourhood
of $e$ such that for every $x$ in it, every \m{k\in\HT{T}\cap\cgp{\cL}} and every vector $y$ with \m{k\,y=x},}
\[\mH{\cL}\bigl(\cgp{\cL}\bigr)\cdot\sum_{\theta\in\Theta_0}\bigl(\rePart\,v_\theta^{*}Rv_\theta\bigr)^{+}\ \le\ \lenergy{Z}{R}{\omega}{y}.\]
\end{enumerate}}
\leanline{\lean{lintegralenergyletopKernel}\sep \lean{eventuallysumpacketsleenergy}}

Theorem~A is the decay of the finite sieve in the frequency, Theorem~B its bottom value, and Theorems C and~D
are what the top floor of the sieve needs: the last form is an integral of a positive kernel, and that kernel
dominates the sum over the packets of rational terms that the main note of the unit (Note~F, \emph{Translates of
Poincaré series fill the rows}) compares with Siegel's formula.

\sect{The setting}{s:setting}
\begin{conventions}\phantomsection\label{c:conventions}
The natural numbers are \m{0,1,2,\dots}, and a difference of natural numbers is read as $0$ when it would be
negative. $A$ is a \emph{semilocal field}: a locally compact topological ring with a Haar measure \m{\vol} and a
compact open subring $O_A$ of volume~$1$, in this proof always a finite product of non-archimedean local fields.
It carries an involution \m{x\mapsto\bar x} and an \emph{integral place structure}: $O_A$ is stable under the
involution; \m{O_A^{\times}} is the group of the units of~$O_A$; $\psi_A$ is a unitary additive character of $A$
which is \emph{self-dual}, that is, \m{\psi_A(ra)=1} for all \m{r\in O_A} exactly if \m{a\in O_A}; for a unit $u$
of~$A$, \m{\modA(u)} is the modulus, \m{\vol(uX)=\modA(u)\vol(X)}, and \m{\gix{u}=\modA(u)^{-1}} is the
\emph{generalised index}, which for \m{u\in O_A} is the index of $uO_A$ in~$O_A$, so that then \m{\gix{u}\ge1}; for a
unit~$\xi$, \m{\den{\xi}} is the index of \m{O_A\cap\xi^{-1}O_A} in~$O_A$. The place structure has further data (a
trace-one element, a measure on the skew elements) that enter only through the inputs; its Plancherel constant
is~$1$. A \emph{shape} is a triple \m{\fs=(\varrho,\kappa,k_0)} of a natural number, a finite set and an element
of it; \m{\rh{\fs}=\varrho}, \m{\kk{\fs}=\#\kappa}, \m{\nn{\fs}=\rh{\fs}+1}. $T$ is an \emph{$\fs$-form} over~$A$
whose split Gram matrix is a \emph{hermitian unit} (hermitian and invertible); $\iota$ is a finite set, $m$ a
natural number, and $\cL$ denotes a level of~$T$ in \m{\Levels{\iota}{T}{m+1}}, in the sense recalled below. A
\emph{sequence of units} is a map \m{p\colon\NN\to A^{\times}}, \m{a\mapsto p_a}; the sequences of units carry the
product topology, and $1$ is the constant sequence. For a finite set $X$, \m{\ell^2(X)} is \m{\CC^X} with the inner
product \m{\langle v,w\rangle=\sum_x\bar v_xw_x}, conjugate-linear in the first variable; $X^*$ is the conjugate
transpose of a matrix. Integrals of functions with values in \m{[0,\infty]} are Lebesgue integrals; for measures
$\mu$, $\nu$ on a group, \m{\mu*\nu} is the image of \m{\mu\times\nu} under the multiplication. Elements written
$x$, $y$, $u$, $k$ are elements of \m{\GT{T}}, and $g$ is a unit of~$A$, where nothing else is said. This note
is local, and some letters have in it the meanings given here, not those of the overview of the write-up
(those in brackets): $\sigma$, $\tau$ are elements of a group (a field automorphism, an infinite place), $s$ is a
natural number (the complex variable), $q$ a sequence of units (a rational number), $\theta$ an index and
$\Theta_0$ an index set (the fit family, the chain of lifts), $\psi_A$ an additive character (a column),
\m{\zetaA{s}} the ideal zeta function of~$O_A$ (Riemann's function), and \m{S(I)} a set of positions (a set of
columns).
\end{conventions}

\subhead{Forms and groups.} These notions belong to the background library. The index set \m{\Ix{\fs}}
consists of \m{e_0,\dots,e_{\varrho-1}}, of the elements of~$\kappa$, of~$z$, and of
\m{e^*_0,\dots,e^*_{\varrho-1}}. The \emph{frame} is \m{\mathfrak f_a=e_a} for \m{a<\varrho} and \m{\mathfrak
f_{\varrho}=z}. An $\fs$-form $T$ is a hermitian matrix indexed by \m{\kappa\sqcup\{z\}}; its \emph{split Gram
matrix} \m{\JT{T}}, indexed by \m{\Ix{\fs}}, pairs $e_a$ with $e^*_a$ to~$1$ and is $T$ on the middle;
\m{\GT{T}} is the unitary group of \m{\JT{T}} over~$A$, with the form \m{h(v,v')=v^*\JT{T}v'}, so that the
coordinate of a vector $w$ at $e^*_a$ is \m{h(e_a,w)}. For \m{x\in\GT{T}},
\m{\ZT{T}{x}=(h(\mathfrak f_a,x\,\mathfrak f_{a'}))_{a,a'\le\varrho}} is its \emph{frame Gram matrix}; its corner
entry is \m{\ZT{T}{x}_{\varrho\varrho}=h(z,xz)}. The \emph{level} of an index is $a$ for $e_a$, $\varrho$ for the
middle indices, and \m{2\varrho-a} for $e^*_a$. In \m{\GT{T}}: \m{\UT{T}} is the \emph{Bessel unipotent
subgroup}, the elements that are block upper triangular for the levels with identity diagonal blocks;
\m{\HT{T}} is the subgroup of the elements that fix all $e_a$, all $e^*_a$ and~$z$; \m{\Hl{l}{T}}, for
\m{l\le\varrho}, the subgroup of those that fix $e_a$ and $e^*_a$ for \m{a<l}, and \m{\Hl{\varrho+1}{T}=\HT{T}};
\m{\Ql{l}{T}} is the group of the elements that are unitriangular on the first $l$ floors; and
\m{\ST{T}=\UT{T}\HT{T}} is the \emph{Bessel subgroup}. For \m{u\in\GT{T}}, \m{\sd{u}=\sum_{a<\varrho}h(e^*_a,u\,\mathfrak
f_{a+1})} is the sum of the entries of $u$ just above the diagonal of its frame block. For a sequence $p$ of
units and \m{l\le\varrho}, \m{\weyl{l}{p}\in\GT{T}} is the \emph{Weyl element} of length~$l$: it sends $e_a$ to
\m{p_ae^*_a} and $e^*_a$ to \m{\bar p_a^{-1}e_a} for \m{a<l} and fixes the other basis vectors. For a unit $g$
of~$A$, \m{\uLev{g}{T}} is the group of the \m{k\in\GT{T}} that have entries in~$O_A$, as has $k^{-1}$, and for
which \m{g^{-1}(k-1)} has entries in~$O_A$; so \m{\uLev{1}{T}} is the group of the \emph{integral isometries}.

\subhead{Levels and the tower.} These notions belong to the unit on the real place (the notes \emph{The tower
method at a place}); they are recalled with their arguments. A \emph{level}
\m{\cL\in\Levels{\iota}{T}{m+1}} consists of a hyperbolic basis of the middle form with \m{m+1} further
hyperbolic pairs and an orthonormal part indexed by~$\iota$; a unit \m{\gL{\cL}} of $A$ with
\m{\overline{\gL{\cL}}=\gL{\cL}}; s-finite measures \m{\flm{\cL}{\varrho}}, carried by \m{\Hl{\varrho}{T}}, and
\m{\mH{\cL}\ne0}, carried by \m{\HT{T}} and invariant under right translation by its elements; an open set of vectors,
the \emph{ball}; a compact set \m{\BoxL{\cL}\subseteq\GT{T}}, the \emph{box}; and a measurable set
\m{\Posn{\cL}\subseteq\GT{T}}, the \emph{positions}. Its \emph{congruence subgroup} is
\m{\cgp{\cL}=\uLev{\gL{\cL}}{T}}. The library provides measures $\eta_l$ on \m{\GT{T}}, the \emph{Heisenberg
measures} of the floors; \m{\etab{j}} is the convolution of the first $j$ of them, and \m{\etab{\varrho+2}} is
carried by \m{\UT{T}}. The \emph{floor measures} are \m{\flm{\cL}{l}} for \m{l<\varrho} (the law of the inverse
Weyl element of floor~$l$, convolved with~$\eta_l$), \m{\flm{\cL}{\varrho}}, and \m{\flm{\cL}{\varrho+1}=\mH{\cL}}.
For natural numbers $l$, $i$ the \emph{stack measure} is
\[\stk{\cL}{l}{i}=(\eta_{l+i-1}*\flm{\cL}{l+i-1})*\dots*(\eta_l*\flm{\cL}{l}),\]
with \m{\gM{\cL}=\stk{\cL}{0}{\varrho+2}}, the \emph{group measure}, and
\m{\gMu{\cL}=\flm{\cL}{\varrho+1}*\dots*\flm{\cL}{0}}. The \emph{quadric measure} \m{\quadric{\cL}{\varrho}} is the
image of \m{\flm{\cL}{\varrho}} under \m{\sigma\mapsto\sigma^{-1}z}, a measure on the vectors \m{e\in A^{\Ix{\fs}}} with
\m{h(e,e)=h(z,z)}. For a unit $g$ and \m{l<\varrho} the \emph{floor modulus} is
\m{\fmod{\cL}{l}{g}=\modA(g)^{\kk{\fs}+1+2(\varrho-l-1)+1/2}}, and for sequences $p$, $q$ of units and
\m{j\le\varrho} the \emph{head density} is
\[\hdens{\cL}{j}{p}{q}=\prod_{a<j}\modA(p_aq_a)^{(\kk{\fs}+1+2(\varrho-a-1))/2}.\]
The \emph{volume exponent} is \m{\vexp{\cL}=\#\iota-(3\nn{\fs}+m+2)}. For \m{a<\varrho} and \m{x,y\in\GT{T}} the set of
\emph{floor values} \m{\Val{\cL}{a}{x}{y}\subseteq A} consists of the coordinates \m{h(e_a,w)} of the vectors $w$
whose coordinates at $e^*_b$, \m{b<a}, vanish and for which \m{x^{-1}w} and
\m{\gL{\cL}^{-1}(x^{-1}w-(xy)^{-1}w)} lie in the ball. A point $x$ is \emph{near at~$j$} if \m{\xi x\in\BoxL{\cL}}
for some \m{\xi\in\Ql{j}{T}}.

For a function \m{k\colon\GT{T}\times\GT{T}\to\CC} and \m{j\le\varrho+1} put
\begin{align*}
\Av{\cL}{j}{k}{g}{y}&=\int k(\sigma g,\sigma y)\,d\bigl(\flm{\cL}{\varrho+1}*\dots*\flm{\cL}{j}\bigr)(\sigma),\\
\Pv{\cL}{j}{k}{g}{y}&=\int\bigl|\Av{\cL}{j+1}{k}{\sigma g}{\sigma y}\bigr|\,d\flm{\cL}{j}(\sigma),
\end{align*}
here with \m{g,y\in\GT{T}}. For a finite set $X$, functions \m{W,W'\colon\GT{T}\to\ell^2(X)} and \m{d\in\GT{T}}, write
$k_{W,W'}^{d}(x,y)=\langle W(x),W'(dy)\rangle$. The \emph{pairing norm} is
$\PN{\cL}{W}{W'}=\int\!\!\int|\langle W(x),W'(y)\rangle|^2\,d\gMu{\cL}(y)\,d\gMu{\cL}(x)$. For \m{j\le\varrho},
sequences $p$, $q$ of units, \m{x,g\in\GT{T}} and \m{\omega\colon\GT{T}\to\CC}, with \m{d=\weyl{j}{q}^{-1}\weyl{j}{p}}, put
\begin{align*}
\Phi_j(x,g)=\int\!\!\int&\omega(\sigma\tau^{-1})\cdot\Av{\cL}{j+1}{k^{1}_{W,W}}{\sigma g}{\sigma gx}\\
&\cdot\overline{\Av{\cL}{j+1}{k^{d}_{W,W}}{\tau\weyl{j}{p}^{-1}g}{\tau\weyl{j}{p}^{-1}gx}}\,d\flm{\cL}{j}(\tau)\,d\flm{\cL}{j}(\sigma).
\end{align*}
Then the \emph{floor form}, its \emph{modulus} at $x$ and its \emph{majorant} at $x$ are
\begin{align*}
\fform{\cL}{j}{p}{q}{W}{\omega}&=\hdens{\cL}{j}{p}{q}\cdot\int\!\!\int\Phi_j(x,g)\,d\gM{\cL}(x)\,d\stk{\cL}{0}{j}(g),\\
\fmodu{\cL}{j}{p}{q}{W}{x}{\omega}&=\int|\Phi_j(x,g)|\,d\stk{\cL}{0}{j}(g),\\
\fmaj{\cL}{j}{p}{q}{W}{W'}{x}&=\int\Pv{\cL}{j}{k^{1}_{W,W}}{g}{gx}\\
&\qquad\cdot\Pv{\cL}{j}{k^{d}_{W,W'}}{\weyl{j}{p}^{-1}g}{\weyl{j}{p}^{-1}gx}\,d\stk{\cL}{0}{j}(g).
\end{align*}
A function \m{W\colon\GT{T}\to\ell^2(X)} is a \emph{cut Bessel function for~$\cL$} if it is measurable and bounded,
\m{W(ux)=\psi_A\bigl(\gL{\cL}^{-1}\sd{u}\bigr)W(x)} for \m{u\in\UT{T}}, and \m{W(x)\ne0} implies \m{ux\in\BoxL{\cL}}
for some \m{u\in\UT{T}}. For a measure $\mu$ on \m{\GT{T}}, such a $W$ and \m{\gamma\in\GT{T}}, the \emph{Bessel
correlation} \m{\bcorr{\mu}{W}{\gamma}\in\ell^2(X\times X)} has the coordinates
\m{\int W(y)_a\,\overline{W(\gamma^{-1}y)_b}\,d\mu(y)}.

A \emph{gauge} \m{\Gamma\in\Gauges{\cL}} provides, for every natural number~$i$, a compact set of the elements
of $A$ \emph{of size at most~$i$}, written \m{\szl{\Gamma}{x}{i}}, increasing with~$i$, and a constant
\m{\Lam{\Gamma}\in[2,\infty)}; a vector or a matrix has size at most $i$ if its entries have. A sequence $p$ of
units is a \emph{head of length $j$ for~$\Gamma$}, \m{\Hd{\Gamma}{j}{p}}, if \m{(p_0\cdots p_{l-1})^{-1}} has size at
most~$1$ for all \m{l\le j}. The \emph{self-norm product}
\m{\snp{\Gamma}{j}{p}{q}{W}{W'}{x}\in[0,\infty]} of a gauge is the quantity of that name of the unit on the real
place; it enters here only through the inputs.

\subhead{Standard levels.} From the note \emph{Poincaré series of cut Bessel functions} of this unit:
\m{\IsStd{\cL}} (\emph{$\cL$ is a standard level}) means that $\cL$ is the level of the lattice
\m{O_A^{\Ix{\fs}}} with a bound for an ideal zeta function; what this note uses of it is Input~\ref{i:3}. For
a unit $\alpha$ of $A$ with \m{\alpha\in O_A} and \m{\bar\alpha=\alpha}, \m{\stdg{\cL}{\alpha}\in\Gauges{\cL}} is the
\emph{gauge of depth~$\alpha$} of a standard level: $x$ has size at most $i$ exactly if \m{\alpha^ix\in O_A}, and
\m{\Lam{\stdg{\cL}{\alpha}}=2/\modA(\alpha)}. For a natural number $K$, \m{j\le\varrho} and a sequence $p$ of units, the
\emph{head weight} is \m{\fhw{K}{j}{p}=\prod_{l<j}\gix{p_0\cdots p_l}^{K/2}} if all the partial products
\m{p_0\cdots p_l}, \m{l<j}, lie in~$O_A$, and \m{\fhw{K}{j}{p}=0} otherwise; \m{\zetaA{s}=\sum_I[O_A:I]^{-s}} is
the ideal zeta function of~$O_A$, the sum over the ideals of finite index.

\begin{setting}[a level]\label{s:L}
$A$ is a semilocal field with an involution and an integral place structure, $\fs$ is a shape, $T$ is an
$\fs$-form over $A$ whose split Gram matrix \m{\JT{T}} is a hermitian unit, $\iota$ is a finite set, $m$ a natural
number, and \m{\cL\in\Levels{\iota}{T}{m+1}} is a level.
\end{setting}

\begin{definition}[the finite Bessel function]\label{d:phase}
Let $A$ be a semilocal field with an involution and an integral place structure, $\fs$ a shape and $T$ an
$\fs$-form over $A$ whose split Gram matrix is a hermitian unit. For a unit $g$ of $A$ and \m{z\in\GT{T}},
\[\phs{g}{z}=\begin{cases}\psi_A\bigl(g^{-1}\sd{u}\bigr)&\text{if }u^{-1}z\in\uLev{g}{T}\text{ for some }u\in\UT{T},\\
0&\text{if there is no such }u,\end{cases}\]
where in the first case $u$ is a chosen such element; by Lemma~\ref{l:phase}(c) the value does not depend on the
choice.
\end{definition}
\leanline{\lean{phase}}

\begin{definition}[intermediate groups and the finite Bessel vector]\label{d:int}
In Setting~\ref{s:L}:
\begin{enumerate}[label=\textup{(\alph*)},leftmargin=2.3em,itemsep=2pt]
\item \m{\IntG{\cL}} is the set of the subgroups $Z$ of \m{\GT{T}} with \m{\cgp{\cL}\le Z\le\uLev{1}{T}}, the
\emph{intermediate groups} of~$\cL$.
\item The \emph{unipotent mass} is \m{\umass{\cL}=\prod_{l<\varrho}\fmod{\cL}{l}{\gL{\cL}}\in[0,\infty]}.
\item For a subgroup $Z$ of \m{\GT{T}},
\[\fzf{\cL}{Z}=\int\Bigl|\int_Z\phs{\gL{\cL}}{yxy^{-1}}\,d\gM{\cL}(y)\Bigr|^2d\gM{\cL}(x)\ \in[0,\infty].\]
\item For \m{Z\in\IntG{\cL}}, \m{\Cos{Z}} is the set of the cosets \m{\cgp{\cL}a}, \m{a\in Z}; it is finite,
because \m{\cgp{\cL}} is open and $Z$ lies in the compact group \m{\uLev{1}{T}}. Each coset $c$ has a chosen
representative \m{a_c\in Z}.
\item For \m{Z\in\IntG{\cL}} the \emph{mass ratio} is
$\mr{Z}=\gM{\cL}(Z)\cdot\gM{\cL}(\cgp{\cL})\cdot\umass{\cL}^{-2}\in[0,\infty]$.
\item For \m{Z\in\IntG{\cL}} and \m{y\in\GT{T}} the \emph{finite Bessel vector} \m{\bvec{Z}{y}\in\ell^2(\Cos{Z})}
has the coordinate \m{\phs{\gL{\cL}}{y\,a_c^{-1}}} at the coset~$c$. We write
\m{\pk{Z}(x,y)=\langle\bvec{Z}{x},\bvec{Z}{y}\rangle} for its kernel.
\end{enumerate}
\end{definition}
\leanline{\lean{IntermediateGroup}\sep \lean{unipotentMass}\sep \lean{floorZeroForm}\sep \lean{Cosets}\sep
\lean{instFintypeCosets} (the finiteness, an instance without a name)\sep \lean{massRatio}\sep \lean{besselVector}}

\begin{setting}[a standard level with an intermediate group]\label{s:Z}
In Setting~\ref{s:L}, with $A$, $\fs$, $T$, $\iota$, $m$ and $\cL$ as there: $\cL$ is a standard level,
\m{\IsStd{\cL}}, and \m{Z\in\IntG{\cL}} is an intermediate group of~$\cL$.
\end{setting}

\begin{definition}[correlations, the top kernel, the energy]\label{d:top}
In Setting~\ref{s:Z}, write \m{\varrho=\rh{\fs}}.
\begin{enumerate}[label=\textup{(\alph*)},leftmargin=2.3em,itemsep=2pt]
\item For a sequence $p$ of units of $A$ and \m{l\in\GT{T}}, the \emph{correlation} is
\[\corr{Z}{p}{l}=\hdens{\cL}{\varrho}{p}{1}\cdot\bcorr{\gM{\cL}}{\bvec{Z}{\cdot}}{\weyl{\varrho}{p}\,l^{-1}}\ \in\ \ell^2(\Cos{Z}\times\Cos{Z}).\]
\item For a set \m{L\subseteq A} and sequences $p$, $q$ of units of~$A$, the \emph{top kernel} is
\[\topK{L}{Z}{p}{q}=\int_{\{l:\ h(z,l^{-1}z)-p_\varrho\in L\}}\bigl\langle\corr{Z}{p}{l},\corr{Z}{q}{l}\bigr\rangle\,d\stk{\cL}{\varrho}{2}(l)\ \in\CC.\]
\item For a finite set $\mathcal C$, a complex matrix $R$ indexed by~$\mathcal C$, a family
\m{\omega=(\omega_i)_{i\in\mathcal C}} of sequences of units of $A$ and a vector \m{e\in A^{\Ix{\fs}}}, the \emph{energy} is
\[\lenergy{Z}{R}{\omega}{e}=\inf_{l}\ \int\sum_{x\in\Cos{Z}\times\Cos{Z}}\bigl(\rePart\,v^{*}Rv\bigr)^{+}\,d\mH{\cL}(h)\ \in[0,\infty],\]
where the infimum is over the \m{l\in\Hl{\varrho}{T}} with \m{l^{-1}z=e} (it is $\infty$ if there is none), $v$ is
the vector with the coordinates \m{v_i=\corr{Z}{\omega_i}{hl}_x}, the index $i$ running over~$\mathcal C$, and
\m{r^{+}=\max(r,0)}. (The index, called $p$ in the formal text, is bound inside the definition; see
§\ref{s:notes}, item~1.)
\end{enumerate}
\end{definition}
\leanline{\lean{correlation}\sep \lean{topKernel}\sep \lean{energya}}

\sect{What is taken for granted}{s:inputs}
The following statements are used and not proved here. Inputs 1 and~5 belong to the background library,
Inputs 2 and~4 to the unit on the real place (the notes \emph{The tower method at a place}), Input~3 to the note
\emph{Poincaré series of cut Bessel functions} of this unit. They are stated from the dossier's readings of
their quoted Lean statements; the statements of Inputs 2 to~4 I also read in the Lean text. No reference to the
literature is claimed for any of them: they are statements about the objects of this proof.

\statement{est}{\phantomsection\label{i:1}\textbf{Input 1} (the Bessel subgroup and the congruence subgroups;
background library).\enspace Let $A$, $\fs$, $T$ be as in the conventions and $g$ a unit of~$A$.
\begin{enumerate}[label=\textup{(\alph*)},leftmargin=2.3em,itemsep=1pt]
\item \m{\UT{T}\subseteq\ST{T}}, and \m{\ST{T}} normalises \m{\UT{T}}. For \m{s,s'\in\ST{T}}:
\m{\sd{ss'}=\sd{s}+\sd{s'}}; and \m{\sd{h}=0} for \m{h\in\HT{T}}.
\item \m{\uLev{g}{T}} is an open compact subgroup of \m{\GT{T}}, contained in \m{\uLev{1}{T}} and normalised by it;
for \m{k\in\uLev{g}{T}}, \m{g^{-1}\sd{k}\in O_A}.
\item If \m{\JT{T}} and its inverse have entries in~$O_A$ and \m{g\in O_A}, then every \m{x\in\GT{T}} for which
\m{g^{-1}(x-1)} has entries in $O_A$ lies in \m{\uLev{g}{T}}.
\item For a sequence $p$ of units and \m{l\le\varrho}, \m{\weyl{l}{p}} commutes with the elements of \m{\Hl{l}{T}}.
\m{\HT{T}\subseteq\Hl{\varrho}{T}}; if \m{\sigma,\tau\in\Hl{\varrho}{T}} and \m{\sigma^{-1}z=\tau^{-1}z}, then
\m{\sigma\tau^{-1}\in\HT{T}}; and \m{\ZT{T}{hxh'}_{\varrho\varrho}=\ZT{T}{x}_{\varrho\varrho}} for \m{h,h'\in\HT{T}} and
\m{x\in\GT{T}}.
\item For \m{u\in\UT{T}}, \m{a<\varrho} and a vector $w$ whose coordinates at $e^*_b$, \m{b<a}, vanish, the
coordinates of $uw$ and of $w$ at $e^*_a$ agree.
\end{enumerate}}

\statement{est}{\phantomsection\label{i:2}\textbf{Input 2} (the tower; unit on the real place).\enspace In
Setting~\ref{s:L}, write \m{\varrho=\rh{\fs}}.
\begin{enumerate}[label=\textup{(\alph*)},leftmargin=2.3em,itemsep=1pt]
\item \emph{(Measures.)} \m{\gM{\cL}=\etab{\varrho+2}*\gMu{\cL}}; \m{\gM{\cL}=\stk{\cL}{\varrho}{2}*\stk{\cL}{0}{\varrho}};
\m{\stk{\cL}{\varrho}{2}=\mH{\cL}*\flm{\cL}{\varrho}}, and this measure is invariant under right translation by the
elements of \m{\Hl{\varrho}{T}}; \m{\stk{\cL}{l}{i}} is carried by \m{\Hl{l}{T}}, and \m{\stk{\cL}{0}{0}} is the unit
mass at~$1$.
\item \emph{(Pairing kernels.)} Let $W$ be a cut Bessel function for $\cL$, \m{k=k^{1}_{W,W}}, and suppose that
$\cL$ has a gauge. Then \m{|\Av{\cL}{j}{k}{g}{y}|\le\Pv{\cL}{j}{k}{g}{y}} for \m{j\le\varrho} and \m{g,y\in\GT{T}}; there
is \m{c<\infty} with \m{\int|k(\sigma g,\sigma y)|\,d\stk{\cL}{\varrho}{2}(\sigma)\le c} for all $g$,~$y$; and for every
function~$k'$, if \m{\Pv{\cL}{j}{k'}{g}{y}\ne0} then \m{k'(\sigma g,\sigma y)\ne0} for some \m{\sigma\in\Hl{j}{T}}.
\item \emph{(Forms.)} For \m{j\le\varrho}, sequences $p$, $q$ of units, $W$ and $\omega$:
\m{|\fform{\cL}{j}{p}{q}{W}{\omega}|\le\hdens{\cL}{j}{p}{q}\cdot\int\fmodu{\cL}{j}{p}{q}{W}{x}{\omega}\,d\gM{\cL}(x)}; and if $W$
is a cut Bessel function for $\cL$ and \m{|\omega|\le1}, then
\m{\fmodu{\cL}{j}{p}{q}{W}{x}{\omega}\le\fmaj{\cL}{j}{p}{q}{W}{W}{x}} for all~$x$.
\item \emph{(Near points.)} If $W$ is a cut Bessel function for $\cL$ and \m{W(x)\ne0}, then $x$ is near at
every~$j$; if $x$ is near at $j$ and \m{\sigma\in\Hl{j}{T}}, then \m{\sigma x} is near at~$j$. If $\Gamma$ is a gauge
of~$\cL$, \m{j\le\varrho}, $p$ is a sequence of units, and
$x$ and \m{\weyl{j}{p}^{-1}x} are near at~$j$, then \m{\szl{\Gamma}{p_0\cdots p_{l-1}}{9\nn{\fs}}} for all \m{l\le j}.
\item \emph{(Correlations.)} Let $\Gamma$ be a gauge of~$\cL$, $W$ a continuous cut Bessel function for $\cL$,
$p$, $q$ heads of length $\varrho$ for~$\Gamma$, \m{x\in\Hl{\varrho}{T}} and $\mu$ a Haar measure of \m{\GT{T}}. Then
\begin{align*}
&\bigl\langle\bcorr{\mu}{W}{\weyl{\varrho}{p}x},\bcorr{\mu}{W}{\weyl{\varrho}{q}x}\bigr\rangle\\
&\quad=\int\!\!\int\langle W(y_1),W(y_2)\rangle\,\overline{\bigl\langle W((\weyl{\varrho}{p}x)^{-1}y_1),W((\weyl{\varrho}{q}x)^{-1}y_2)\bigr\rangle}\,d\mu\,d\mu,
\end{align*}
the integral over \m{(y_1,y_2)}. If \m{W(y)\ne0} and \m{W((\weyl{\varrho}{p}x)^{-1}y)\ne0}, then \m{\szl{\Gamma}{y}{60\nn{\fs}^2}}; and for every $i$ the
set of the \m{y\in\GT{T}} with \m{\szl{\Gamma}{y}{i}} is compact.
\end{enumerate}}

\statement{est}{\phantomsection\label{i:3}\textbf{Input 3} (standard levels; note \emph{Poincaré series of cut
Bessel functions} of this unit).\enspace In Setting~\ref{s:L}, suppose \m{\IsStd{\cL}}, and write
\m{\varrho=\rh{\fs}}, \m{n=\nn{\fs}}.
\begin{enumerate}[label=\textup{(\alph*)},leftmargin=2.3em,itemsep=1pt]
\item \m{\JT{T}} and its inverse have entries in~$O_A$; \m{\gL{\cL}\in O_A}; \m{\BoxL{\cL}=\uLev{1}{T}}; the ball is
\m{O_A^{\Ix{\fs}}}; \m{\Posn{\cL}=\{y^{-1}uy:\ y\in\GT{T},\ u\in\UT{T}\}\cdot\cgp{\cL}}; and there is a vector $u_0$
with \m{h(u_0,w)=w_z} for all vectors~$w$ and with $z$-coordinate~$0$.
\item \m{\zetaA{m}^{n}\le1+2^{-n}}, hence \m{\zetaA{m}\le2}; \m{m<9n}; \m{13n\le\#\iota}; and
\m{\kk{\fs}=\#\iota+2m+3}.
\item \m{\gM{\cL}} is a Haar measure of \m{\GT{T}}, invariant under left and right translations; and
\m{\vol(O_A^{\times})^{i+1}\le\stk{\cL}{l}{i+2}(\uLev{1}{T})} whenever \m{l+i=\varrho}.
\item \emph{(Unfolding.)} For a unit $g$ of $A$ with \m{g\in O_A} and \m{\bar g=g}, an s-finite measure $\theta$
on \m{\GT{T}} and \m{c\in\GT{T}}:
\begin{align*}
&(\etab{\varrho+2}*\theta)\bigl(\{y:\ yc\in\uLev{g}{T}\}\bigr)\\
&\qquad=\prod_{l<\varrho}\fmod{\cL}{l}{g}\cdot\theta\bigl(\{y:\ uyc\in\uLev{g}{T}\text{ for some }u\in\UT{T}\}\bigr).
\end{align*}
\item \emph{(Gauges.)} For a unit \m{\alpha\in O_A} with \m{\bar\alpha=\alpha}, \m{\stdg{\cL}{\alpha}} is a gauge of~$\cL$,
as described in §\ref{s:setting}; in every gauge $\Gamma$ of~$\cL$ the elements of the box and of the ball have
size at most~$1$, the sizes add under multiplication, and \m{\modA(u)\le\Lam{\Gamma}^{i}} for a unit $u$ with
\m{\szl{\Gamma}{u}{i}}.
\end{enumerate}}

\statement{est}{\phantomsection\label{i:4}\textbf{Input 4} (the estimates of the tower method; unit on the real
place).\enspace In Setting~\ref{s:L}, suppose \m{\IsStd{\cL}}; let $W$ be a cut Bessel function for~$\cL$,
\m{k=k^{1}_{W,W}}, \m{j\le\rh{\fs}}, $p$, $q$ sequences of units, $\Gamma$ a gauge of~$\cL$, and write
\m{n=\nn{\fs}}, \m{\varrho=\rh{\fs}}.
\begin{enumerate}[label=\textup{(\alph*)},leftmargin=2.3em,itemsep=1pt]
\item \emph{(Positions.)} Let $q$ be a head of length $j$ for~$\Gamma$, $\Gamma'$ a second gauge, $S$ a set of
positions $x$ with \m{\szl{\Gamma'}{\gL{\cL}^{-1}(x-1)}{1}}, and \m{f\colon\GT{T}\to[0,\infty]}, \m{\epsilon\in[0,\infty]} with
\m{f(x)^2\le\epsilon^2\cdot\snp{\Gamma}{j}{p}{q}{W}{W}{x}} on~$S$. Then
\[\hdens{\cL}{j}{p}{q}\int_Sf\,d\gM{\cL}\ \le\ \epsilon\cdot\umass{\cL}\cdot\PN{\cL}{W}{W}\cdot
\Lam{\Gamma}^{N_1}\cdot\Lam{\Gamma'}^{N_2},\]
where \m{N_1=1500000n^3+1} and \m{N_2=50000n^2}.
\item \emph{(The majorant.)} Let $p$ be a head of length $j$ for~$\Gamma$, $x$ a position and
\m{\phi_0\in[0,\infty]}. Suppose that, if \m{\phi_0<1}, for every $y$ with \m{\langle W(y),W(yx)\rangle\ne0}
such that \m{\weyl{j}{p}^{-1}y} and \m{\weyl{j}{q}^{-1}yx} are near at~$j$ there is $a$ with \m{j\le a<\varrho} and
\m{\vol(\Val{\cL}{a}{y}{x})\le\phi_0}. Then
\m{\fmaj{\cL}{j}{p}{q}{W}{W}{x}^2\le\phi_0^{\vexp{\cL}}\cdot V_0^2\cdot\snp{\Gamma}{j}{p}{q}{W}{W}{x}} with
\m{V_0=2^{99\cdot13\,n}}, the floor-mass constant of the gauge of depth~$1$.
\item \emph{(Frequencies.)} Let $p$ be a head of length $j$ for~$\Gamma$ and $x$ a position. If $\eta$, $\xi$ are
units with \m{\psi_A(\xi c)=\psi_A(\eta c)} for all $c$ with \m{\szl{\Gamma}{c}{200n^2}}, and
\m{\omega(y)=\psi_A(\xi\cdot\ZT{T}{y}_{jj})}, then
\m{\fmodu{\cL}{j}{p}{q}{W}{x}{\omega}^2\le\gix{\eta}^{\#\iota}\cdot\snp{\Gamma}{j}{p}{q}{W}{W}{x}}. If \m{j=i+1} and
\m{\xi=\gL{\cL}^{-1}(p_i^{-1}-q_i^{-1})} is a unit, then
\m{\fmodu{\cL}{j}{p}{q}{W}{x}{1}^2\le\gix{\xi}^{\#\iota}\cdot\snp{\Gamma}{j}{p}{q}{W}{W}{x}}.
\item \emph{(Integrality from floor values.)} Let $p$, $q$ be heads of length $j$ for~$\Gamma$ and $\Gamma''$ a
second gauge such that, if \m{j\ne0}, \m{\szl{\Gamma''}{\gL{\cL}^{-1}c}{1}} for every $c$ with \m{\szl{\Gamma}{c}{10n}}.
Let \m{x,y\in\GT{T}} be such that: there are \m{u,v\in\UT{T}} with \m{ux,\,vxy\in\BoxL{\cL}} and
\m{\szl{\Gamma''}{\gL{\cL}^{-1}(vxy-ux)}{2}}; \m{\weyl{j}{p}^{-1}x} and \m{\weyl{j}{q}^{-1}xy} are near at~$j$; and for
every $a$ with \m{j\le a<\varrho} there is a unit \m{c\in\Val{\cL}{a}{x}{y}} with \m{\szl{\Gamma''}{c^{-1}}{1}}. Then
\m{\szl{\Gamma''}{\gL{\cL}^{-1}(y-1)}{99n}}.
\item \emph{(The floor $0$.)} \m{\hdens{\cL}{0}{p}{q}=1},
\m{\fmaj{\cL}{0}{1}{1}{W}{W}{x}=\Pv{\cL}{0}{k}{1}{x}^2}, and there is a number \m{\Sigma_0(x)\in[0,\infty]} with
\m{\snp{\Gamma}{0}{1}{1}{W}{W}{x}=\Lam{\Gamma}^{3\cdot10^6n^3}\cdot\Sigma_0(x)^2} such that, for a position $x$ and
\m{\varepsilon\in[0,\infty]}: if for every \m{\sigma\in\GT{T}} with \m{k(\sigma,\sigma x)\ne0} there is \m{a<\varrho} with
\m{\vol(\Val{\cL}{a}{\sigma}{x})\le\varepsilon}, then
\m{\Pv{\cL}{0}{k}{1}{x}^2\le\varepsilon^{\vexp{\cL}}\cdot\Lam{\Gamma}^{10^4n\varrho}\cdot\Sigma_0(x)}.
\end{enumerate}}

\statement{est}{\phantomsection\label{i:5}\textbf{Input 5} (ideals, valuations, Eichler elements; background
library).\enspace Let $A$ be as in the conventions.
\begin{enumerate}[label=\textup{(\alph*)},leftmargin=2.3em,itemsep=1pt]
\item \emph{(Summation over ideals.)} Let $s$ be a natural number with \m{\zetaA{s}\le2}, $\rho$ a measure on a
measurable space, $S$ a subset of it, $f$ a function with values in \m{[0,\infty]}, \m{V\in[0,\infty]}, and
\m{I\mapsto S(I)} a family of subsets indexed by the ideals of~$O_A$. If for almost every \m{x\in S} there is an
ideal $I$ of finite index with \m{x\in S(I)} and \m{[O_A:I]^{s}\cdot f(x)\cdot\rho(S(I))\le V}, then
\m{\int_Sf\,d\rho\le\zetaA{s}\cdot V}.
\item \emph{(Content.)} For a finite family \m{(d_i)} of elements of~$A$ one of which is~$1$, there is a unit
$c$ of $A$ with \m{c\in O_A} such that \m{\{\varepsilon\in A:\ \varepsilon d_i\in O_A\text{ for all }i\}=cO_A}; this is an
ideal of $O_A$ of index \m{\gix{c}}, and \m{\bar c\in O_A}.
\item \emph{(Greatest common divisors.)} For units $\xi$, $\beta$ of $A$ there are a unit $\eta$ and
\m{x,c\in O_A} with \m{\eta=\xi+c\beta}, \m{\xi=x\eta} and \m{\beta\in\eta O_A}.
\item \emph{(Primitive matrices.)} If a matrix $X$ over~$A$, with columns indexed by a finite set with a
weight function, has a left inverse, there is an invertible matrix $b$, block upper triangular for the
weights, such that $Xb$ has entries in $O_A$ and a left inverse with entries in~$O_A$.
\item \emph{(Volumes.)} For a unit~$u$, \m{\vol(\{x:\ ux\in O_A\})=\modA(u)^{-1}}; for a unit \m{u\in O_A} the ideal
\m{uO_A} has index \m{\gix{u}}; and \m{\vol(O_A\cap\xi^{-1}O_A)=\den{\xi}^{-1}}. Almost every element of $A$ is a
unit.
\item \emph{(Eichler elements.)} In Setting~\ref{s:L}, let $u_0$ be a vector with \m{h(u_0,w)=w_z} for all $w$
and with $z$-coordinate~$0$. For every vector $v$ there is \m{\mathrm E(v)\in\GT{T}}, given by a formula that is
polynomial in the coordinates of $v$ and in the ring inverse of~$v_z$, with \m{\mathrm E(z)=1}, such that
\m{\mathrm E(v)\in\Hl{\varrho}{T}} if the coordinates of $v$ at all $e_a$ and $e^*_a$ vanish, and
\m{\mathrm E(v)^{-1}z=v} if \m{h(v,v)=h(z,z)} and $v_z$ is a unit.
\end{enumerate}}

\sect{The finite Bessel function and its vector}{s:phase}
\begin{lemma}[the finite Bessel function]\label{l:phase}
Let $A$ be a semilocal field with an involution and an integral place structure, $\fs$ a shape, $T$ an
$\fs$-form over~$A$, $g$ a unit of~$A$, and \m{x,y,u,k\in\GT{T}}.
\begin{enumerate}[label=\textup{(\alph*)},leftmargin=2.3em,itemsep=1pt]
\item If \m{\phs{g}{y}\ne0}, there is \m{u'\in\UT{T}} with \m{u'^{-1}y\in\uLev{g}{T}}.
\item \m{|\phs{g}{y}|\le1}.
\item If \m{u\in\UT{T}} and \m{u^{-1}y\in\uLev{g}{T}}, then \m{\phs{g}{y}=\psi_A(g^{-1}\sd{u})}.
\item If \m{x\in\ST{T}}, \m{u\in\UT{T}} and \m{u^{-1}x\in\uLev{g}{T}}, then
\m{\phs{g}{xy}=\psi_A(g^{-1}\sd{u})\cdot\phs{g}{y}}.
\item If \m{k\in\uLev{g}{T}}, then \m{\phs{g}{yk}=\phs{g}{y}}.
\item If \m{u\in\UT{T}}, then \m{\phs{g}{uy}=\psi_A(g^{-1}\sd{u})\cdot\phs{g}{y}}.
\item The function \m{z\mapsto\phs{g}{z}} on \m{\GT{T}} is locally constant, hence continuous.
\item \m{\uLev{g}{T}\le\uLev{1}{T}}.
\end{enumerate}
\end{lemma}
\begin{proof}
(a) and (b) are read off Definition~\ref{d:phase}, as $\psi_A$ is unitary. (h) is part of the definition of
\m{\uLev{g}{T}}.

(c) Let $u_0$ be the chosen element, so that \m{u_0^{-1}y\in\uLev{g}{T}}. Then
\m{u_0^{-1}u=(u_0^{-1}y)(u^{-1}y)^{-1}} lies in \m{\uLev{g}{T}}, so \m{g^{-1}\sd{u_0^{-1}u}\in O_A} by
Input~\ref{i:1}(b), and \m{\psi_A(g^{-1}\sd{u_0^{-1}u})=1} because $\psi_A$ is self-dual. As
\m{\sd{u}=\sd{u_0}+\sd{u_0^{-1}u}} by Input~\ref{i:1}(a), the two values of the character agree.

(d) Suppose first that \m{u_0^{-1}y\in\uLev{g}{T}} for some \m{u_0\in\UT{T}}. Then \m{v=xu_0x^{-1}u} lies in
\m{\UT{T}}, because $x$ normalises it, and $v^{-1}xy=(u^{-1}x)(u_0^{-1}y)\in\uLev{g}{T}$. By (c),
$\phs{g}{xy}=\psi_A(g^{-1}\sd{v})$, and $\sd{v}=\sd{x}+\sd{u_0}+\sd{x^{-1}}+\sd{u}=\sd{u_0}+\sd{u}$ by additivity.
With $\phs{g}{y}=\psi_A(g^{-1}\sd{u_0})$ this is the claim. If there is no such $u_0$, then \m{\phs{g}{y}=0}, and
\m{\phs{g}{xy}=0} as well: from \m{v^{-1}xy\in\uLev{g}{T}} with \m{v\in\UT{T}} one gets \m{u_0=x^{-1}vu^{-1}x\in\UT{T}}
with $u_0^{-1}y=(u^{-1}x)^{-1}(v^{-1}xy)\in\uLev{g}{T}$.

(e) The condition \m{u^{-1}yk\in\uLev{g}{T}} is the condition \m{u^{-1}y\in\uLev{g}{T}}; so both sides vanish
together, and otherwise they are equal by~(c). (f) is (d) for \m{x=u}. (g) By (e) the function is constant on
the cosets of the open subgroup \m{\uLev{g}{T}} (Input~\ref{i:1}(b)).
\end{proof}
\leanline{\lean{existsbesselUofphasenezero}\sep \lean{normphaseleone}\sep \lean{phaseeq}\sep
\lean{phasebesselSmul}\sep \lean{phasemulofmemunitaryLevel}\sep \lean{phasebesselUmul}\sep
\lean{continuousphase}\sep \lean{unitaryLevelleunitaryLevelone}}

\begin{lemma}[intermediate groups]\label{l:int}
In Setting~\ref{s:L}:
\begin{enumerate}[label=\textup{(\alph*)},leftmargin=2.3em,itemsep=1pt]
\item \m{\umass{\cL}\ne0} and \m{\umass{\cL}\ne\infty}.
\item Every \m{Z\in\IntG{\cL}} is an open subgroup of \m{\GT{T}}.
\item For \m{Z\in\IntG{\cL}} and \m{y\in Z} there is a coset \m{d\in\Cos{Z}} such that, for every
\m{c\in\Cos{Z}}, \m{y\,a_c^{-1}\in\cgp{\cL}} holds exactly if \m{c=d}.
\item For every unit $u$ of~$A$, \m{\modA(u)^{-1}=\bigl(\sqrt{\gix{u}}\bigr)^2}.
\end{enumerate}
Suppose moreover that $\cL$ is a standard level, \m{\IsStd{\cL}}.
\begin{enumerate}[label=\textup{(\alph*)},leftmargin=2.3em,itemsep=1pt,resume]
\item If $S$ is a subgroup of \m{\GT{T}} contained in \m{\uLev{1}{T}}, then \m{\gM{\cL}(S)\ne\infty}.
\item For every \m{Z\in\IntG{\cL}}: \m{\mr{Z}\ne0} and \m{\mr{Z}\ne\infty}.
\end{enumerate}
\end{lemma}
\begin{proof}
(a) Each floor modulus is a product of powers of \m{\modA(\gL{\cL})}, the modulus of a unit, which is neither $0$
nor~$\infty$. (b) $Z$ contains the open subgroup \m{\cgp{\cL}} (Input~\ref{i:1}(b)). (c) Take for $d$ the coset
of~$y$: \m{y\,a_c^{-1}\in\cgp{\cL}} says that $y$ and $a_c$ lie in the same coset. (d) is the definition of the
generalised index. (e) \m{\gM{\cL}} is a Haar measure (Input~\ref{i:3}(c)) and \m{\uLev{1}{T}} is compact. (f) $Z$
and \m{\cgp{\cL}} are open, hence of positive measure, and of finite measure by~(e); with (a) this gives the claim.
\end{proof}
\leanline{\lean{unipotentMassnezeroandnetop}\sep \lean{isOpen}\sep \lean{existscoset}\sep
\lean{invmdeqsqrtgenIndexsq}\sep \lean{groupMeasurenetop}\sep \lean{massRationezeroandnetop}}

\begin{lemma}[non-orthogonal Bessel vectors]\label{l:inner}
In Setting~\ref{s:L}, let \m{Z\in\IntG{\cL}} and \m{x,y\in\GT{T}} with \m{\langle\bvec{Z}{x},\bvec{Z}{y}\rangle\ne0}.
Then there are \m{u,v\in\UT{T}} with \m{ux\in\uLev{1}{T}} and \m{(ux)^{-1}(vy)\in\cgp{\cL}}.
\end{lemma}
\begin{proof}
For some coset $c$ both $\phs{\gL{\cL}}{x\,a_c^{-1}}$ and $\phs{\gL{\cL}}{y\,a_c^{-1}}$ are not zero. By
Lemma~\ref{l:phase}(a) there are \m{u,v\in\UT{T}} such that $k_1=ux\,a_c^{-1}$ and $k_2=vy\,a_c^{-1}$ lie in
\m{\cgp{\cL}}.
Then \m{ux=k_1a_c\in\uLev{1}{T}}, as \m{a_c\in Z\le\uLev{1}{T}}, and
$(ux)^{-1}(vy)=a_c^{-1}(k_1^{-1}k_2)a_c\in\cgp{\cL}$, because \m{\uLev{1}{T}} normalises \m{\cgp{\cL}}
(Input~\ref{i:1}(b)).
\end{proof}
\leanline{\lean{existsofinnerbesselVectornezero}}

\begin{lemma}[the form on the congruence subgroup]\label{l:cong}
In Setting~\ref{s:Z}, with $x$, $y$ the variables of integration and \m{g=\gL{\cL}},
\[\int_{\cgp{\cL}}\Bigl|\int_Z\phs{g}{yxy^{-1}}\,d\gM{\cL}(y)\Bigr|^2d\gM{\cL}(x)=\gM{\cL}(Z)^2\cdot\gM{\cL}\bigl(\cgp{\cL}\bigr).\]
\end{lemma}
\begin{proof}
For \m{x\in\cgp{\cL}} and \m{y\in Z\le\uLev{1}{T}} the conjugate \m{yxy^{-1}} lies in \m{\cgp{\cL}}, so
\m{\phs{g}{yxy^{-1}}=\psi_A(g^{-1}\sd{1})=1} by Lemma~\ref{l:phase}(c) with \m{u=1}. The inner integral is
\m{\gM{\cL}(Z)}, which is finite by Lemma~\ref{l:int}(e).
\end{proof}
\leanline{\lean{floorZeroFormoncongGroup}}

\begin{lemma}[the finite Bessel vector is cut]\label{l:cut}
In Setting~\ref{s:Z}, the function \m{y\mapsto\bvec{Z}{y}} is a continuous cut Bessel function for~$\cL$.
\end{lemma}
\begin{proof}
Its coordinates are continuous by Lemma~\ref{l:phase}(g) and of modulus at most~$1$, so its norm is at most
\m{\sqrt{\#\Cos{Z}}}. The rule \m{\bvec{Z}{uy}=\psi_A(\gL{\cL}^{-1}\sd{u})\bvec{Z}{y}} for \m{u\in\UT{T}} is
Lemma~\ref{l:phase}(f) in each coordinate. If \m{\bvec{Z}{x}\ne0}, then \m{\langle\bvec{Z}{x},\bvec{Z}{x}\rangle\ne0},
and Lemma~\ref{l:inner} gives \m{u\in\UT{T}} with \m{ux\in\uLev{1}{T}=\BoxL{\cL}} (Input~\ref{i:3}(a)).
\end{proof}
\leanline{\lean{isCutbesselVector}}

\begin{proposition}[unfolding along the unipotent subgroup]\label{p:unfold}
In Setting~\ref{s:Z}, for all \m{y,x\in\GT{T}}, with \m{g=\gL{\cL}} and $u$ the variable of integration,
\[\umass{\cL}\cdot\bigl\langle\bvec{Z}{y},\bvec{Z}{yx}\bigr\rangle=\int\mathbf 1_Z(uy)\cdot\phs{g}{uyx(uy)^{-1}}\,d\etab{\varrho+2}(u).\]
\end{proposition}
\begin{proof}
The measure \m{\etab{\varrho+2}} is carried by \m{\UT{T}}; let \m{u\in\UT{T}}. We claim that the integrand at $u$
equals
\[\sum_{c\in\Cos{Z}}\mathbf 1\bigl[uy\,a_c^{-1}\in\cgp{\cL}\bigr]\cdot\overline{\bvec{Z}{y}_c}\cdot\bvec{Z}{yx}_c .\]
If \m{uy\notin Z}, both vanish, since \m{uy\,a_c^{-1}\in\cgp{\cL}} would give \m{uy\in\cgp{\cL}a_c\subseteq Z}. If
\m{uy\in Z}, by Lemma~\ref{l:int}(c) exactly one coset $d$ has \m{k=uy\,a_d^{-1}\in\cgp{\cL}}, and the sum is
$\overline{\phs{g}{y\,a_d^{-1}}}\cdot\phs{g}{yx\,a_d^{-1}}$. Now \m{y\,a_d^{-1}=u^{-1}k}, so
$\phs{g}{y\,a_d^{-1}}=\psi_A(g^{-1}\sd{u^{-1}})=\overline{\psi_A(g^{-1}\sd{u})}$ by Lemma~\ref{l:phase}(c) and
Input~\ref{i:1}(a); and $uyx(uy)^{-1}=u\cdot(yx\,a_d^{-1})\cdot k^{-1}$, so
$\phs{g}{uyx(uy)^{-1}}=\psi_A(g^{-1}\sd{u})\cdot\phs{g}{yx\,a_d^{-1}}$ by Lemma~\ref{l:phase}(e),~(f). This proves
the claim. Integrating it, the coset $c$ contributes
$\etab{\varrho+2}(\{u:\ uy\,a_c^{-1}\in\cgp{\cL}\})\cdot\overline{\bvec{Z}{y}_c}\,\bvec{Z}{yx}_c$. By the unfolding of
Input~\ref{i:3}(d) for the unit mass at~$1$, this measure is \m{\umass{\cL}} if \m{u'y\,a_c^{-1}\in\cgp{\cL}} for some
\m{u'\in\UT{T}}, which is the case when \m{\bvec{Z}{y}_c\ne0} (Lemma~\ref{l:phase}(a)); and when
\m{\bvec{Z}{y}_c=0} the term vanishes. The sum over $c$ is $\umass{\cL}\langle\bvec{Z}{y},\bvec{Z}{yx}\rangle$.
\end{proof}
\leanline{\lean{unipotentMassmulinnerbesselVector}}

\begin{proposition}[the pairing norm]\label{p:pn}
In Setting~\ref{s:Z},
$\PN{\cL}{\bvec{Z}{\cdot}}{\bvec{Z}{\cdot}}\cdot\umass{\cL}\le\mr{Z}$.
\end{proposition}
\begin{proof}
Write \m{W=\bvec{Z}{\cdot}}, \m{\mathsf v=\umass{\cL}}, and \m{\nu(x)=\etab{\varrho+2}(\{u:\ ux\in Z\})} for
\m{x\in\GT{T}}. By Proposition~\ref{p:unfold} (for the pair \m{x}, \m{x^{-1}y}) and Lemma~\ref{l:phase}(b),
\begin{equation}\label{e:beta}\mathsf v\cdot|\langle W(x),W(y)\rangle|\le\nu(x)\qquad(x,y\in\GT{T}).\end{equation}
Moreover \m{\nu(x)\le1}: as \m{Z\le\uLev{1}{T}}, Input~\ref{i:3}(d) for the unit~$1$ and the unit mass at $1$
bounds \m{\nu(x)} by \m{\prod_{l<\varrho}\fmod{\cL}{l}{1}=1}. Since \m{\gM{\cL}=\etab{\varrho+2}*\gMu{\cL}}
(Input~\ref{i:2}(a)), \m{\gM{\cL}(Z)=\int\nu(x)\,d\gMu{\cL}(x)}. So it suffices to show, for every~$x$,
\begin{equation}\label{e:pn}\mathsf v^{3}\int|\langle W(x),W(y)\rangle|^2\,d\gMu{\cL}(y)\ \le\ \nu(x)\cdot\gM{\cL}\bigl(\cgp{\cL}\bigr),\end{equation}
and to integrate over~$x$. If \m{W(x)=0} this is clear. Otherwise, as $W$ is cut (Lemma~\ref{l:cut}), there is
\m{u_0\in\UT{T}} with \m{b=u_0x\in\uLev{1}{T}}. By the right invariance of \m{\gM{\cL}} and Input~\ref{i:3}(d) for
\m{\theta=\gMu{\cL}} and \m{c=b^{-1}},
\begin{gather*}
\gM{\cL}\bigl(\cgp{\cL}\bigr)=\gM{\cL}\bigl(\{y:\ yb^{-1}\in\cgp{\cL}\}\bigr)=\mathsf v\cdot\gMu{\cL}(Y_x),\\
Y_x=\{y:\ uyb^{-1}\in\cgp{\cL}\text{ for some }u\in\UT{T}\}.
\end{gather*}
For \m{y\in Y_x}, \eqref{e:beta} and \m{\nu(x)\le1} give \m{\mathsf v^2|\langle W(x),W(y)\rangle|^2\le\nu(x)^2\le\nu(x)}.
For \m{y\notin Y_x} the inner product vanishes: otherwise Lemma~\ref{l:inner} gives \m{u',v'\in\UT{T}} with
\m{u'x\in\uLev{1}{T}} and \m{v'y\in u'x\,\cgp{\cL}}; then
\m{v'yb^{-1}\in u'u_0^{-1}\cdot b\,\cgp{\cL}\,b^{-1}=u'u_0^{-1}\cgp{\cL}}, since $b$ normalises \m{\cgp{\cL}}, and
\m{u=u_0u'^{-1}v'} shows \m{y\in Y_x}. Hence the left side of \eqref{e:pn} is at most
\m{\mathsf v\cdot\nu(x)\cdot\gMu{\cL}(Y_x)=\nu(x)\gM{\cL}(\cgp{\cL})}. Integrating \eqref{e:pn} over $x$ gives
\m{\mathsf v^3\cdot\PN{\cL}{W}{W}\le\gM{\cL}(Z)\,\gM{\cL}(\cgp{\cL})}, which is the claim after division
by~$\mathsf v^2$ (Lemma~\ref{l:int}(a)).
\end{proof}
\leanline{\lean{pairingNormmullemassRatio}}

\begin{lemma}[the inner integral as an average]\label{l:avg}
In Setting~\ref{s:Z}, for every \m{x\in\GT{T}}, with $y$ the variable of integration and \m{g=\gL{\cL}},
\[\Bigl|\int_Z\phs{g}{yxy^{-1}}\,d\gM{\cL}(y)\Bigr|=\bigl|\Av{\cL}{0}{\pk{Z}}{1}{x}\bigr|\cdot\umass{\cL}.\]
\end{lemma}
\begin{proof}
The integrand is bounded on $Z$, which has finite measure. By $\gM{\cL}=\etab{\varrho+2}*\gMu{\cL}$ and Fubini's
theorem the integral is
\[\int\!\!\int\mathbf 1_Z(u\sigma)\,\phs{g}{u\sigma x(u\sigma)^{-1}}\,d\etab{\varrho+2}(u)\,d\gMu{\cL}(\sigma),\]
and by Proposition~\ref{p:unfold} the inner integral is $\umass{\cL}\cdot\pk{Z}(\sigma,\sigma x)$. The integral of
\m{\pk{Z}(\sigma,\sigma x)} over $\sigma$ is \m{\Av{\cL}{0}{\pk{Z}}{1}{x}} by definition.
\end{proof}
\leanline{\lean{enormintegralphaseequpperPairing}}

\begin{lemma}[a lower bound for the mass of the integral group]\label{l:mass}
In Setting~\ref{s:L}, suppose \m{\IsStd{\cL}}, and let \m{B\in[0,\infty]} with \m{2\le B} and
\m{\vol(O_A\setminus O_A^{\times})\le B^{-1}}. Then \m{2^{-\nn{\fs}}\le\gM{\cL}(\uLev{1}{T})}.
\end{lemma}
\begin{proof}
By Input~\ref{i:3}(c) for \m{l=0}, \m{i=\varrho}: $\vol(O_A^{\times})^{\varrho+1}\le\gM{\cL}(\uLev{1}{T})$; and
$\vol(O_A^{\times})=1-\vol(O_A\setminus O_A^{\times})\ge1-B^{-1}\ge\tfrac12$.
\end{proof}
\leanline{\lean{invtwopowlegroupMeasure}}

\sect{The shell bound, the frequency bound and the bottom value}{s:shells}
In this section \m{n=\nn{\fs}} and \m{\varrho=\rh{\fs}}; in the proofs \m{g=\gL{\cL}} and \m{W=\bvec{Z}{\cdot}}. For a
standard level and the gauge \m{\stdg{\cL}{\alpha}}, ``size at most $i$'' means \m{\alpha^ix\in O_A}; for
\m{\alpha=1} it means integrality.

\begin{lemma}[floor values]\label{l:val}
In Setting~\ref{s:Z}, let \m{a<\rh{\fs}} and \m{x,z,y\in\GT{T}} with \m{\langle\bvec{Z}{x},\bvec{Z}{z}\rangle\ne0}. Then
there is a unit $c$ of $A$ with \m{c\in O_A}, \m{c\in\Val{\cL}{a}{x}{y}} and
\m{\vol\bigl(\Val{\cL}{a}{x}{y}\bigr)\le\modA(c)}.
\end{lemma}
\begin{proof}
Let $X$ be the matrix with the rows of \m{x^{-1}} above the rows of \m{g^{-1}(x^{-1}-(xy)^{-1})}; it has a left
inverse. As the ball is \m{O_A^{\Ix{\fs}}} (Input~\ref{i:3}(a)), \m{\Val{\cL}{a}{x}{y}} is the set of the coordinates
$w_{e^*_a}$ of the vectors $w$ with $Xw$ integral and \m{w_{e^*_b}=0} for \m{b<a}. By Input~\ref{i:5}(d), with the
levels of the indices as weights, there is an invertible matrix $b$, block upper triangular for the levels, such
that $Xb$ is integral with an integral left inverse; so $Xw$ is integral exactly if \m{b^{-1}w} is. Put
\m{i=e^*_a}. Its level \m{2\varrho-a} is shared by no other index, and the indices of higher level are the
$e^*_b$ with \m{b<a}. As $b$ and $b^{-1}$ are block upper triangular, the column $w^0$ of $b$ at $i$ vanishes at
these indices, and \m{(b^{-1}w)_i=(b^{-1})_{ii}\,w_i} for every $w$ that vanishes at them. For \m{w=w^0} this gives
\m{(b^{-1})_{ii}\,b_{ii}=1}; let \m{c=b_{ii}}, a unit. Since $Xw^0$ is a column of~$Xb$, \m{c=w^0_i\in\Val{\cL}{a}{x}{y}}.
Every element of \m{\Val{\cL}{a}{x}{y}} is a $w_i$ with \m{b^{-1}w} integral and $w$ vanishing at the $e^*_b$,
\m{b<a}; then \m{c^{-1}w_i=(b^{-1}w)_i\in O_A}. So \m{\Val{\cL}{a}{x}{y}\subseteq cO_A}, a set of volume \m{\modA(c)}.
Finally, by Lemma~\ref{l:inner} there is \m{u\in\UT{T}} with \m{ux\in\uLev{1}{T}}; \m{x^{-1}w^0} is integral, hence so
is \m{uw^0=(ux)(x^{-1}w^0)}, and its coordinate at $e^*_a$ is $w^0_i=c$ by Input~\ref{i:1}(e). So \m{c\in O_A}.
\end{proof}
\leanline{\lean{existsunitmemfloorValues}}

\begin{lemma}[non-orthogonal vectors are close]\label{l:close}
In Setting~\ref{s:Z}, let $\Gamma$ be a gauge of~$\cL$ and \m{x,y\in\GT{T}} with
\m{\langle\bvec{Z}{x},\bvec{Z}{y}\rangle\ne0}. Then there are \m{u,v\in\UT{T}} with \m{ux\in\BoxL{\cL}},
\m{vy\in\BoxL{\cL}} and \m{\szl{\Gamma}{\gL{\cL}^{-1}(vy-ux)}{2}}.
\end{lemma}
\begin{proof}
Take $u$, $v$ from Lemma~\ref{l:inner}. Then $ux\in\uLev{1}{T}=\BoxL{\cL}$, and $k=(ux)^{-1}(vy)$ lies in
\m{\cgp{\cL}}, so that $vy=(ux)k\in\uLev{1}{T}$. Further $g^{-1}(vy-ux)=(ux)\cdot g^{-1}(k-1)$ is a product of integral matrices; its
columns lie in the ball, whose elements have size at most~$1$ in every gauge (Input~\ref{i:3}(e)).
\end{proof}
\leanline{\lean{existscloseofinnernezero}}

\begin{proposition}[a gauge adapted to two heads]\label{p:depth}
In Setting~\ref{s:Z}, let \m{j\le\rh{\fs}}, let $K$ be a natural number and $p$, $q$ sequences of units of~$A$,
and suppose that \m{\fmaj{\cL}{j}{p}{q}{\bvec{Z}{\cdot}}{\bvec{Z}{\cdot}}{x}\ne0} for some \m{x\in\GT{T}}. Then there is
a unit $\alpha$ of $A$ with \m{\alpha\in O_A} and \m{\bar\alpha=\alpha} such that, for the gauge
\m{\Gamma=\stdg{\cL}{\alpha}}:
\begin{enumerate}[label=\textup{(\roman*)},leftmargin=2.3em,itemsep=1pt]
\item \m{\Hd{\Gamma}{j}{p}} and \m{\Hd{\Gamma}{j}{q}}, and if \m{j\ne0} then \m{\szl{\Gamma}{\gL{\cL}^{-1}}{1}};
\item \m{\szl{\Gamma}{p_i^{-1}-q_i^{-1}}{1}} for all \m{i<j};
\item \m{\Lam{\Gamma}^{E}\le\bigl(2/\modA(\gL{\cL})\bigr)^{K}\cdot\fhw{K}{j}{p}\cdot\fhw{K}{j}{q}} for every natural
number $E$ with \m{4E\le K}.
\end{enumerate}
\end{proposition}
\begin{proof}
Write \m{P_l=p_0\cdots p_{l-1}} and \m{Q_l=q_0\cdots q_{l-1}}. Since the majorant at $x$ is not zero, there is a point
$y$ in the support of \m{\stk{\cL}{0}{j}} at which both pivot moduli of its integrand are not zero, hence
(Input~\ref{i:2}(b)) \m{\sigma,\tau\in\Hl{j}{T}} with \m{\langle W(\sigma y),W(\sigma yx)\rangle\ne0} and
\m{\langle W(\tau\weyl{j}{p}^{-1}y),W(\weyl{j}{q}^{-1}\weyl{j}{p}\tau\weyl{j}{p}^{-1}yx)\rangle\ne0}. As $\tau$ commutes
with the Weyl elements (Input~\ref{i:1}(d)), $W$ does not vanish at \m{\sigma y}, \m{\tau\weyl{j}{p}^{-1}y},
\m{\sigma yx} and \m{\tau\weyl{j}{q}^{-1}yx}. By Input~\ref{i:2}(d), with Lemma~\ref{l:cut}, the points $y$ and
\m{\weyl{j}{p}^{-1}y} are near at~$j$, and so are $yx$ and \m{\weyl{j}{q}^{-1}yx}; hence, for the gauge of depth~$1$,
all \m{P_l} and \m{Q_l} with \m{l\le j} lie in~$O_A$. Put \m{\beta=\prod_{l\le j}P_l} and \m{\gamma=\prod_{l\le j}Q_l}.
Then \m{\beta\in O_A}, \m{\beta/P_l\in O_A} for \m{l\le j}, \m{\beta/p_i=P_i\cdot(\beta/P_{i+1})\in O_A} for \m{i<j}, and
\m{\fhw{K}{j}{p}=\gix{\beta}^{K/2}} by the multiplicativity of the index; likewise for $\gamma$ and~$q$. Let
\[\alpha=g\cdot\overline{\beta\gamma}\cdot\beta\gamma,\]
a unit in $O_A$ with \m{\bar\alpha=\alpha} (Input~\ref{i:3}(a)). Then \m{\alpha P_l^{-1}}, \m{\alpha Q_l^{-1}} (\m{l\le j}),
\m{\alpha g^{-1}} and \m{\alpha(p_i^{-1}-q_i^{-1})} (\m{i<j}) lie in~$O_A$, which is (i) and~(ii). For (iii),
$\Lam{\Gamma}=2\gix{\alpha}=2\gix{g}\,\gix{\beta}^2\gix{\gamma}^2$, and the three indices are at least~$1$; so for
\m{4E\le K}, $\Lam{\Gamma}^E\le2^K\gix{g}^K(\gix{\beta}\gix{\gamma})^{K/2}$.
\end{proof}
\leanline{\lean{existsgaugedepthofnezero}}

\begin{proposition}[the shell bound]\label{p:shells}
In Setting~\ref{s:Z}, let \m{j\le\rh{\fs}}, let $t$ be a natural number with \m{t+10^{8}\nn{\fs}^4\le\kk{\fs}}, let
$\alpha$ be a unit of $A$ with \m{\alpha\in O_A} and \m{\bar\alpha=\alpha}, and let $\Gamma$ be the gauge
\m{\Gamma=\stdg{\cL}{\alpha}}. Let $p$, $q$ be sequences of units of $A$ with \m{\Hd{\Gamma}{j}{p}}, \m{\Hd{\Gamma}{j}{q}}
and, if \m{j\ne0}, \m{\szl{\Gamma}{\gL{\cL}^{-1}}{1}}. Let \m{f\colon\GT{T}\to[0,\infty]} and \m{e\in[0,\infty]} be such
that \m{f(x)\le\fmaj{\cL}{j}{p}{q}{\bvec{Z}{\cdot}}{\bvec{Z}{\cdot}}{x}} for all~$x$ and such that, if \m{e<1}, then
\m{f(x)^2\le\bigl(e^{\vexp{\cL}}\bigr)^2\cdot\snp{\Gamma}{j}{p}{q}{\bvec{Z}{\cdot}}{\bvec{Z}{\cdot}}{x}} for all
\m{x\in\Posn{\cL}}. Then
\[\hdens{\cL}{j}{p}{q}\cdot\int f\,d\gM{\cL}\ \le\ \min(1,e)^{t}\cdot\mr{Z}\cdot\Lam{\Gamma}^{10^{9}\nn{\fs}^4}.\]
\end{proposition}
\begin{proof}
Write \m{\Lambda=\Lam{\Gamma}}, \m{s=\vexp{\cL}}. By Input~\ref{i:3}(b) and the hypothesis on~$t$,
\begin{equation}\label{e:s}s=\kk{\fs}-3m-3n-5\ \ge\ 8\cdot10^{7}n^4+t .\end{equation}

\emph{Step 1: $f$ lives on the positions.} If \m{f(x)\ne0}, the majorant at $x$ is not zero; as in the proof of
Proposition~\ref{p:depth} there is $y'$ with $\langle W(y'),W(y'x)\rangle\ne0$, and Lemma~\ref{l:inner} gives
\m{u,v\in\UT{T}} with \m{(uy')^{-1}(vy'x)\in\cgp{\cL}}, that is, $x\in y'^{-1}(v^{-1}u)y'\cdot\cgp{\cL}\subseteq\Posn{\cL}$
(Input~\ref{i:3}(a)).

\emph{Step 2: shells.} For a position $x$ let $I(x)$ be the set of the \m{\varepsilon\in A} that multiply $1$ and
all entries of \m{g^{-1}(x-1)} into~$O_A$. By Input~\ref{i:5}(b), \m{I(x)=cO_A} for a unit \m{c=c(x)\in O_A}, an
ideal of index \m{\gix{c}}. For an ideal $I$ let $S(I)$ be the set of the positions with \m{I(x)=I}. By
Input~\ref{i:5}(a) for the measure \m{f\,d\gM{\cL}}, the constant function \m{\hdens{\cL}{j}{p}{q}} and the
exponent~$m$ (recall \m{\zetaA{m}\le2}), the claim follows, with \m{2\cdot2^{10^5n^2}\le\Lambda^{10^8n^4}}, once we
show for every position $x$, with \m{c=c(x)} and \m{S=S(cO_A)}:
\begin{equation}\label{e:shell}\gix{c}^{m}\cdot\hdens{\cL}{j}{p}{q}\cdot\int_{S}f\,d\gM{\cL}\ \le\ \min(1,e)^{t}\cdot\mr{Z}\cdot\Lambda^{9\cdot10^{8}n^4}\cdot2^{10^{5}n^2}.\end{equation}

\emph{Step 3: one shell.} Let \m{\Gamma'=\stdg{\cL}{\bar cc}}, so \m{\Lam{\Gamma'}=2\gix{c}^2}, and let
\m{\phi\in[0,\infty]} be the number with
\[\phi^{400n^2}=\modA(c)\cdot\Lambda^{(10n+1)\cdot99n}.\]
For \m{x'\in S} the matrix $\bar cc\cdot g^{-1}(x'-1)$ is integral, so $S$ consists of positions $x'$ with
$\szl{\Gamma'}{g^{-1}(x'-1)}{1}$. Put
\[\epsilon=\min(\phi,1,e)^{s}\cdot\Lambda^{9999n}.\]
We show in Step~4 that $f(x')^2\le\epsilon^2\cdot\snp{\Gamma}{j}{p}{q}{W}{W}{x'}$ for all \m{x'\in S}. Then Input~\ref{i:4}(a), with
Lemma~\ref{l:cut} and Proposition~\ref{p:pn}, gives
\[\hdens{\cL}{j}{p}{q}\cdot\int_Sf\,d\gM{\cL}\ \le\ \epsilon\cdot\mr{Z}\cdot\Lambda^{1500000n^3+1}\cdot\bigl(2\gix{c}^2\bigr)^{50000n^2}.\]
By \eqref{e:s},
\begin{align*}\min(\phi,1,e)^{s}&\le\phi^{400n^2\cdot200000n^2}\cdot\min(1,e)^{t}\\
&=\modA(c)^{200000n^2}\cdot\Lambda^{(10n+1)99n\cdot200000n^2}\cdot\min(1,e)^t.\end{align*}
The powers of \m{\gix{c}=\modA(c)^{-1}} are
\m{m+100000n^2-200000n^2\le0}, as \m{m<9n}; they are at most~$1$. The exponent of $\Lambda$ is
\[9999n+1500000n^3+1+(10n+1)\cdot99n\cdot200000n^2\ \le\ 9\cdot10^8n^4.\]
This is \eqref{e:shell}.

\emph{Step 4: the bound on a shell.} Let \m{x'\in S}. If \m{e<\min(\phi,1)}, then \m{e<1} and the hypothesis
gives \m{f(x')^2\le(e^s)^2\,\snp{\Gamma}{j}{p}{q}{W}{W}{x'}\le\epsilon^2\,\snp{\Gamma}{j}{p}{q}{W}{W}{x'}}. Otherwise
\m{\min(\phi,1,e)=\min(\phi,1)}, and by \m{f\le\fmaj{\cL}{j}{p}{q}{W}{W}{\cdot}} and Input~\ref{i:4}(b)
for \m{\phi_0=\min(\phi,1)^2}, with \m{V_0^2\le\Lambda^{2\cdot9999n}}, it suffices to verify the hypothesis of that
input. So let \m{\phi<1}, and let $y$ be such that \m{\langle W(y),W(yx')\rangle\ne0} and \m{\weyl{j}{p}^{-1}y},
\m{\weyl{j}{q}^{-1}yx'} are near at~$j$; suppose, for a contradiction, that
\m{\vol(\Val{\cL}{a}{y}{x'})>\phi^2} for all $a$ with \m{j\le a<\varrho}. By Lemma~\ref{l:val} there are units
\m{d_a\in O_A} with \m{d_a\in\Val{\cL}{a}{y}{x'}} and \m{\vol(\Val{\cL}{a}{y}{x'})\le\modA(d_a)}. Let
\m{C=\prod_{j\le a<\varrho}d_a} and let $\Gamma''$ be the gauge of depth \m{\alpha^{10n+1}\bar CC}. Then
\m{C/d_a\in O_A}, so \m{d_a^{-1}} has size at most $1$ for~$\Gamma''$; if \m{j\ne0} and $c'$ has size at most
\m{10n} for~$\Gamma$, then \m{\alpha^{10n+1}g^{-1}c'\in O_A}, so \m{g^{-1}c'} has size at most $1$ for~$\Gamma''$; and
Lemma~\ref{l:close} provides the two unipotent elements. Input~\ref{i:4}(d), for the pair \m{(y,x')}, shows that
\m{\varepsilon=(\alpha^{10n+1}\bar CC)^{99n}} multiplies \m{g^{-1}(x'-1)} into~$O_A$; it also lies in~$O_A$. So
\m{\varepsilon\in I(x')=cO_A}, whence \m{\modA(\varepsilon)\le\modA(c)}, that is,
\[\modA(C)^{198n}\ \le\ \modA(c)\cdot\modA(\alpha)^{-(10n+1)99n}\ \le\ \modA(c)\cdot\Lambda^{(10n+1)99n}=\phi^{400n^2}.\]
On the other hand \m{\modA(C)=\prod_a\modA(d_a)>\phi^{2(\varrho-j)}}, or \m{\modA(C)=1} if \m{j=\varrho}; as
\m{\phi<1} and \m{396n(\varrho-j)\le400n^2}, this gives \m{\modA(C)^{198n}>\phi^{400n^2}}, a contradiction.
\end{proof}
\leanline{\lean{lintegralshellsle}}

\begin{corollary}[the majorant is integrable]\label{c:finite}
In Setting~\ref{s:Z}, let \m{j\le\rh{\fs}}, let $r$ be a natural number with \m{r+10^{8}\nn{\fs}^4\le\kk{\fs}}, and
let $p$, $q$ be sequences of units of~$A$. Then
\m{\int\fmaj{\cL}{j}{p}{q}{\bvec{Z}{\cdot}}{\bvec{Z}{\cdot}}{x}\,d\gM{\cL}(x)\ne\infty}.
\end{corollary}
\begin{proof}
If the integral is infinite, the majorant is not identically zero; Proposition~\ref{p:depth}, with \m{K=0},
gives a gauge \m{\Gamma=\stdg{\cL}{\alpha}} that satisfies the hypotheses of Proposition~\ref{p:shells}, which we
apply to $f$ the majorant, \m{e=1} and \m{t=r}. The right side is finite (Lemma~\ref{l:int}(f)) and
\m{\hdens{\cL}{j}{p}{q}\ne0}; a contradiction.
\end{proof}
\leanline{\lean{lintegralformMajorantnetop}}

\begin{proof}[Proof of Theorem A]
The function $\omega$ has modulus at most~$1$, so by Input~\ref{i:2}(c)
\begin{equation}\label{e:A}\begin{gathered}|\fform{\cL}{j}{p}{q}{W}{\omega}|\le\hdens{\cL}{j}{p}{q}\cdot\int f\,d\gM{\cL},\\
f(x)=\fmodu{\cL}{j}{p}{q}{W}{x}{\omega}\le\fmaj{\cL}{j}{p}{q}{W}{W}{x}.\end{gathered}\end{equation}
If the majorant vanishes identically, the form is zero. Otherwise let \m{\Gamma=\stdg{\cL}{\alpha}} be the gauge of
Proposition~\ref{p:depth}, \m{\Lambda=\Lam{\Gamma}}, and \m{\beta=\alpha^{200n^2}}.

\emph{A unit $\eta$.} We claim that there are units $\eta$, $w$ of $A$ with
\begin{equation}\label{e:eta}w\in O_A,\quad\xi w\in O_A,\quad\xi\eta^{-1}\in O_A,\quad\min(1,\gix{\eta})\le\modA(w)\cdot\modA(\beta)^{-1},\end{equation}
such that \m{f(x)^2\le\gix{\eta}^{\#\iota}\cdot\snp{\Gamma}{j}{p}{q}{W}{W}{x}} for all positions~$x$. In the second case
of the theorem (\m{j=i+1}, \m{\omega=1}), take \m{\eta=\xi} and \m{w=\beta}: by (i) and (ii) of
Proposition~\ref{p:depth}, $\xi$ has size at most $2$ for~$\Gamma$, so \m{\xi\beta\in O_A}, and the bound is the
second part of Input~\ref{i:4}(c). In the first case, Input~\ref{i:5}(c) gives a unit $\eta$ and \m{x_0,c\in O_A}
with \m{\eta=\xi+c\beta}, \m{\xi=x_0\eta} and \m{\beta\in\eta O_A}; put \m{w=\eta^{-1}\beta}. Then \m{w\in O_A},
\m{\xi w=x_0\beta\in O_A}, \m{\xi\eta^{-1}=x_0}, and \m{\modA(w)\modA(\beta)^{-1}=\gix{\eta}}. For every $c'$ of size at
most \m{200n^2} for~$\Gamma$, \m{(\eta-\xi)c'=c\,\beta c'\in O_A}, so \m{\psi_A(\eta c')=\psi_A(\xi c')}, and the bound is
the first part of Input~\ref{i:4}(c).

\emph{The shell bound.} We apply Proposition~\ref{p:shells} to $f$, the number \m{e=\gix{\eta}^{1/2}} and
\m{t+12} in the place of~$t$; note \m{t+12+10^8n^4\le\kk{\fs}}. Its second hypothesis holds: if \m{e<1}, then
\m{\gix{\eta}^{\#\iota}\le\gix{\eta}^{\vexp{\cL}}=(e^{\vexp{\cL}})^2}, because \m{\vexp{\cL}\le\#\iota}. With \eqref{e:A} we get
\[|\fform{\cL}{j}{p}{q}{W}{\omega}|\ \le\ \min(1,e)^{t}\cdot\min(1,\gix{\eta})^{6}\cdot\mr{Z}\cdot\Lambda^{10^9n^4}.\]
Now $\gix{\eta}\le\gix{\xi\eta^{-1}}\gix{\eta}=\gix{\xi}$, as \m{\xi\eta^{-1}\in O_A}; so \m{\min(1,e)^t\le\gix{\xi}^{t/2}}.
By \eqref{e:eta}, $\min(1,\gix{\eta})^6\le\modA(w)^6\modA(\beta)^{-6}$; here $\modA(\beta)^{-6}\le\Lambda^{1200n^2}$
(Input~\ref{i:3}(e), as \m{\beta^{-1}} has size at most \m{200n^2}), and
$\modA(w)=\vol(wO_A)\le\vol(O_A\cap\xi^{-1}O_A)=\den{\xi}^{-1}$ by Input~\ref{i:5}(e), since \m{wO_A} is contained
in \m{O_A\cap\xi^{-1}O_A}. Finally $\Lambda^{10^9n^4+1200n^2}\le\Lambda^{2\cdot10^9n^4}\le(2/\modA(g))^K\fhw{K}{j}{p}\fhw{K}{j}{q}$
by Proposition~\ref{p:depth}(iii), as \m{8\cdot10^9n^4\le K}. This is the bound of Theorem~A.
\end{proof}

\begin{proof}[Proof of Theorem B]
Let \m{\Gamma=\stdg{\cL}{1}}, so \m{\Lam{\Gamma}=2}, and \m{\epsilon_0=\min\bigl(\vol(O_A\setminus O_A^{\times}),\tfrac12\bigr)}.
Split the outer integral of \m{\fzf{\cL}{Z}} over \m{\cgp{\cL}} and its complement. The first part is
\m{\gM{\cL}(Z)^2\gM{\cL}(\cgp{\cL})} by Lemma~\ref{l:cong}. In the second, by Lemma~\ref{l:avg} and
Input~\ref{i:2}(b) (with Lemma~\ref{l:cut}), the integrand is at most \m{\umass{\cL}^2\cdot f(x)}, where
\[f(x)=\mathbf 1[x\notin\cgp{\cL}]\cdot\Pv{\cL}{0}{\pk{Z}}{1}{x}^2=\mathbf 1[x\notin\cgp{\cL}]\cdot\fmaj{\cL}{0}{1}{1}{W}{W}{x}\]
(Input~\ref{i:4}(e)). We apply Proposition~\ref{p:shells} with \m{j=0}, \m{p=q=1}, \m{e=\epsilon_0} and
\m{t+1+n+10^9n^4} in the place of~$t$; the sequence $1$ is a head of length~$0$, and
\m{t+1+n+10^9n^4+10^8n^4\le\kk{\fs}}. Granting its second hypothesis, and since \m{\hdens{\cL}{0}{1}{1}=1},
\begin{align*}\umass{\cL}^2\int f\,d\gM{\cL}\ &\le\ \umass{\cL}^2\,\mr{Z}\cdot\epsilon_0^{\,t+1}\cdot\epsilon_0^{\,n}\cdot(2\epsilon_0)^{10^9n^4}\\
&\le\ \gM{\cL}(Z)\,\gM{\cL}\bigl(\cgp{\cL}\bigr)\cdot\vol\bigl(O_A\setminus O_A^{\times}\bigr)^{t+1}\cdot2^{-n},\end{align*}
which with the first part is the claim.

It remains to show \m{f(x)^2\le(\epsilon_0^{\vexp{\cL}})^2\cdot\snp{\Gamma}{0}{1}{1}{W}{W}{x}} for a position
\m{x\notin\cgp{\cL}}. By Input~\ref{i:4}(e), and since \m{2\cdot10^4n\varrho\le3\cdot10^6n^3}, this follows once
we show: for every $\sigma$ with \m{\pk{Z}(\sigma,\sigma x)\ne0} there is \m{a<\varrho} with
\m{\vol(\Val{\cL}{a}{\sigma}{x})\le\epsilon_0}. By Lemma~\ref{l:val} there are units \m{u_a\in O_A} with
\m{u_a\in\Val{\cL}{a}{\sigma}{x}} and \m{\vol(\Val{\cL}{a}{\sigma}{x})\le\modA(u_a)}. If some $u_a$ is not a unit
of~$O_A$, then \m{u_aO_A} is a proper ideal contained in \m{O_A\setminus O_A^{\times}}, so
\m{\modA(u_a)=\vol(u_aO_A)\le\epsilon_0}, and $a$ is the floor we want. Otherwise all \m{u_a^{-1}} are integral. Then
Input~\ref{i:4}(d), for \m{\Gamma''=\Gamma}, \m{j=0} and the pair \m{(\sigma,x)}, applies: the unipotent elements come
from Lemma~\ref{l:close}, and every point is near at~$0$, because \m{\Ql{0}{T}=\GT{T}}. It shows that
\m{g^{-1}(x-1)} is integral, so \m{x\in\cgp{\cL}} by Inputs \ref{i:1}(c) and \ref{i:3}(a), against the choice of~$x$.
\end{proof}

\sect{The top floor}{s:top}
Again \m{n=\nn{\fs}}, \m{\varrho=\rh{\fs}}, \m{g=\gL{\cL}} and \m{W=\bvec{Z}{\cdot}}.

\begin{lemma}[one gauge for two sequences]\label{l:heads}
In Setting~\ref{s:L}, suppose \m{\IsStd{\cL}}, and let $p$, $q$ be sequences of units of~$A$. Then there is a
gauge \m{\Gamma\in\Gauges{\cL}} such that all sequences $r$ of units in a neighbourhood of~$p$, and all sequences
$r$ in a neighbourhood of~$q$, are heads of length \m{\rh{\fs}} for~$\Gamma$.
\end{lemma}
\begin{proof}
For \m{l\le\varrho} the map \m{(c,r)\mapsto\bar cc\cdot(r_0\cdots r_{l-1})^{-1}} from pairs of an element of $A$ and a
sequence of units to $A$ is continuous and sends \m{(0,p)} and \m{(0,q)} to \m{0\in O_A}, and $O_A$ is open. So there
are a neighbourhood $N$ of $0$ and neighbourhoods of $p$ and of $q$ on whose products with $N$ all these maps
have values in~$O_A$. Almost every element of $A$ is a unit (Input~\ref{i:5}(e)), so $N$ contains a unit~$c$,
and \m{\Gamma=\stdg{\cL}{\bar cc}} has the property; note \m{\bar cc\in O_A}, the case \m{l=0}.
\end{proof}
\leanline{\lean{existsgaugeisHead}}

\begin{proposition}[the correlations are locally constant]\label{p:const}
In Setting~\ref{s:Z}, let $p$ be a sequence of units of~$A$. Then there is a neighbourhood of \m{(p,1)} in the
product of the space of sequences of units with \m{\GT{T}} such that for every \m{(p',l')} in it with
\m{l'\in\Hl{\varrho}{T}} and every \m{l\in\Hl{\varrho}{T}}: \m{\corr{Z}{p'}{ll'}=\corr{Z}{p}{l}}.
\end{proposition}
\begin{proof}
Let $\Gamma$ be the gauge of Lemma~\ref{l:heads} for $p$ and~$1$, and $C$ the set of the \m{y\in\GT{T}} with
\m{\szl{\Gamma}{y}{60n^2}}, a compact set (Input~\ref{i:2}(e)). The neighbourhood consists of the \m{(p',l')} such that
$p'$ is a head of length $\varrho$ for~$\Gamma$, \m{p_a^{-1}p'_a\in O_A^{\times}} for \m{a<\varrho}, and
\begin{equation}\label{e:k}k_y=y^{-1}\,\weyl{\varrho}{p}\,l'\,\weyl{\varrho}{p'}^{-1}\,y\in\cgp{\cL}\qquad\text{for all }y\in C .\end{equation}
These are neighbourhood conditions: \m{O_A^{\times}} and \m{\cgp{\cL}} are open, the expression in \eqref{e:k} is
continuous in \m{(p',l',y)} and equal to $1$ at \m{(p,1,y)}, and $C$ is compact. Let \m{(p',l')} be such a pair with
\m{l'\in\Hl{\varrho}{T}}, and \m{l\in\Hl{\varrho}{T}}. First, \m{\modA(p'_a)=\modA(p_a)} for \m{a<\varrho}, so
\m{\hdens{\cL}{\varrho}{p'}{1}=\hdens{\cL}{\varrho}{p}{1}}. Secondly, we show for every \m{y\in\GT{T}}
\[W(y)=0\qquad\text{or}\qquad W\bigl(ll'\,\weyl{\varrho}{p'}^{-1}y\bigr)=W\bigl(l\,\weyl{\varrho}{p}^{-1}y\bigr),\]
which gives the equality of the two Bessel correlations. Suppose \m{W(y)\ne0} and one of the two values on the
right is not zero. By Input~\ref{i:2}(e), for the head $p'$ and \m{(ll')^{-1}}, or for the head $p$ and
\m{l^{-1}}, we get \m{y\in C}. Then \m{ll'\weyl{\varrho}{p'}^{-1}y=l\,\weyl{\varrho}{p}^{-1}y\cdot k_y} with
\m{k_y\in\cgp{\cL}}, and for every coset $c$ the element \m{a_ck_ya_c^{-1}} lies in \m{\cgp{\cL}} too, so the
coordinates at $c$ of the two vectors agree by Lemma~\ref{l:phase}(e).
\end{proof}
\leanline{\lean{eventuallycorrelationeq}}

\begin{proof}[Proof of Theorem C]
\emph{(i)} Let $\Gamma$ be the gauge of Lemma~\ref{l:heads} for $p$ and~$q$, so that $p$ and $q$ are heads of
length~$\varrho$. For \m{l\in\Hl{\varrho}{T}}, Input~\ref{i:2}(e) with \m{x=l^{-1}} and \m{\mu=\gM{\cL}}
(Lemma~\ref{l:cut}, Input~\ref{i:3}(c)) gives
\begin{align*}&\bigl\langle\corr{Z}{p}{l},\corr{Z}{q}{l}\bigr\rangle\\
&\quad=\hdens{\cL}{\varrho}{p}{q}\int\!\!\int\langle W(x),W(y)\rangle\\
&\qquad\qquad\cdot\bigl\langle W(l\,\weyl{\varrho}{q}^{-1}y),W(l\,\weyl{\varrho}{p}^{-1}x)\bigr\rangle\,d\gM{\cL}(x)\,d\gM{\cL}(y),\end{align*}
because \m{\hdens{\cL}{\varrho}{p}{1}\hdens{\cL}{\varrho}{q}{1}=\hdens{\cL}{\varrho}{p}{q}}. Let \m{\mathcal A(l,x,y)} be
\m{\varphi_\xi(l^{-1})} times the integrand on the right. It is integrable for
\m{\stk{\cL}{\varrho}{2}\times\gM{\cL}\times\gM{\cL}}: it vanishes unless $x$ and $y$ lie in the compact set of
Input~\ref{i:2}(e), and for fixed $x$, $y$ the integral of its modulus over $l$ is bounded, by the bound for
pairing kernels of Input~\ref{i:2}(b), since
\m{|\langle W(lw_1),W(lw_2)\rangle|=|k^{1}_{W,W}(lw_1,lw_2)|}. So $F_\xi$ is integrable, and
\m{\int F_\xi\,d\stk{\cL}{\varrho}{2}=\hdens{\cL}{\varrho}{p}{q}\int\mathcal A}. We transform \m{\int\mathcal A}. Replace
$y$ by $xx'$ (left invariance of \m{\gM{\cL}}); write \m{x=sy'} with \m{s=h\sigma}, where
\m{(h,\sigma,y')} has the law \m{\mH{\cL}\times\flm{\cL}{\varrho}\times\stk{\cL}{0}{\varrho}} (Input~\ref{i:2}(a)); and
replace $l$ by \m{l's^{-1}}, \m{l'=k\tau} with \m{(k,\tau)} of law \m{\mH{\cL}\times\flm{\cL}{\varrho}}, which is allowed
since \m{s\in\Hl{\varrho}{T}} almost surely and \m{\stk{\cL}{\varrho}{2}} is invariant under right translation by it.
As the Weyl elements commute with \m{\Hl{\varrho}{T}} (Input~\ref{i:1}(d)),
\begin{gather*}l\,\weyl{\varrho}{p}^{-1}x=k\tau\,\weyl{\varrho}{p}^{-1}y'=:z_0,\qquad
l\,\weyl{\varrho}{q}^{-1}xx'=\weyl{\varrho}{q}^{-1}\weyl{\varrho}{p}\,z_0\,x',\\
\varphi_\xi(l^{-1})=\varphi_\xi(h\sigma\tau^{-1}k^{-1})=\varphi_\xi(\sigma\tau^{-1}),\end{gather*}
the last by Input~\ref{i:1}(d). Hence, with \m{d=\weyl{\varrho}{q}^{-1}\weyl{\varrho}{p}},
\[\int\mathcal A=\int\varphi_\xi(\sigma\tau^{-1})\cdot\langle W(h\sigma y'),W(h\sigma y'x')\rangle\cdot
\overline{\bigl\langle W(z_0),W(d\,z_0x')\bigr\rangle},\]
the integral over \m{(k,\tau,h,\sigma,x',y')}. Integrating first over $h$ and over $k$ turns the two inner
products into \m{\Av{\cL}{\varrho+1}{k^{1}_{W,W}}{\sigma y'}{\sigma y'x'}} and the conjugate of
\m{\Av{\cL}{\varrho+1}{k^{d}_{W,W}}{\tau\weyl{\varrho}{p}^{-1}y'}{\tau\weyl{\varrho}{p}^{-1}y'x'}}; integrating then over $\sigma$
and $\tau$ gives \m{\Phi_\varrho(x',y')} for \m{\omega=\varphi_\xi}, and the integral over \m{(x',y')} for
\m{\gM{\cL}\times\stk{\cL}{0}{\varrho}}, times the head density, is the floor form.

\emph{(ii)} Let \m{(p_0,q_0)} be a pair of sequences. By Proposition~\ref{p:const} with \m{l'=1}, for \m{(p,q)}
near \m{(p_0,q_0)} the correlations \m{\corr{Z}{p}{l}=\corr{Z}{p_0}{l}} and \m{\corr{Z}{q}{l}=\corr{Z}{q_0}{l}} agree
for all \m{l\in\Hl{\varrho}{T}}, which carries the measure; and \m{p_\varrho-(p_0)_\varrho\in L} for $p$ near~$p_0$,
because $L$ is open, so the domain of integration is the same, $L$ being a subgroup. The kernel is locally
constant.
\end{proof}

\begin{proof}[Proof of Theorem D]
\emph{(i)} For \m{l\in\GT{T}} and \m{x\in\Cos{Z}\times\Cos{Z}} let \m{v(l,x)} be the vector with the coordinates
\m{\corr{Z}{\omega_i}{l}_x}, \m{i\in\mathcal C}. Then
\[\sum_{i,i'}R_{ii'}\bigl\langle\corr{Z}{\omega_i}{l},\corr{Z}{\omega_{i'}}{l}\bigr\rangle=\sum_xv(l,x)^{*}R\,v(l,x)\ \ge0,\]
as $R$ is positive semidefinite. Each summand of $\mathcal Q$ is an integral of an integrable function
(Theorem~C(i) for \m{\xi=0}), over the same set \m{D=\{l:\ h(z,l^{-1}z)-t\in L\}} because \m{(\omega_i)_\varrho=t}. So
\m{\mathcal Q=\int_D\sum_xv(l,x)^*Rv(l,x)\,d\stk{\cL}{\varrho}{2}(l)} is real and \m{\ge0}, and equals the integral of
\m{G(l)=\sum_x(\rePart\,v(l,x)^*Rv(l,x))^{+}} over~$D$. By \m{\stk{\cL}{\varrho}{2}=\mH{\cL}*\flm{\cL}{\varrho}}, and
since \m{(h\sigma)^{-1}z=\sigma^{-1}z} for \m{h\in\HT{T}},
\begin{align*}\mathcal Q&=\int\mathbf 1\bigl[h(z,\sigma^{-1}z)-t\in L\bigr]\int G(h\sigma)\,d\mH{\cL}(h)\,d\flm{\cL}{\varrho}(\sigma)\\
&\ge\int\mathbf 1\bigl[h(z,\sigma^{-1}z)-t\in L\bigr]\cdot\lenergy{Z}{R}{\omega}{\sigma^{-1}z}\,d\flm{\cL}{\varrho}(\sigma),\end{align*}
because \m{\sigma\in\Hl{\varrho}{T}} almost surely and the energy at \m{\sigma^{-1}z} is an infimum over a set that
contains \m{l=\sigma}. The right side is at least the integral of the theorem: \m{\quadric{\cL}{\varrho}} is the
image of \m{\flm{\cL}{\varrho}} under \m{\sigma\mapsto\sigma^{-1}z}, and \m{h(z,e)=(\JT{T}e)_z}; the integral of a function
for an image measure is at most the integral of its composite, with equality for measurable functions, so no
measurability of the energy is needed.

\emph{(ii)} Fix \m{o\in\Theta_0}; then \m{l_oe=z}. Let $u_0$ be the vector of Input~\ref{i:3}(a) and
\m{\mathrm E} the Eichler elements of Input~\ref{i:5}(f). For a vector $x$ put
\m{m(x)=l_o^{-1}\,\mathrm E(l_ox)\,l_o}. The $z$-coordinate of \m{l_ox} is $1$ at \m{x=e}; as the units are open and
inversion is continuous at~$1$, for $x$ near $e$ this coordinate is a unit and \m{m(x)} is close to
\m{m(e)=1}. By Proposition~\ref{p:const}, applied to each of the finitely many \m{\omega_i}, there is a neighbourhood
$N$ of $e$ such that for \m{x\in N}: the $z$-coordinate of $l_ox$ is a unit, and
\begin{equation}\label{e:m}\corr{Z}{\omega_i}{l\,m(x)}=\corr{Z}{\omega_i}{l}\qquad\text{for all }i\in\mathcal C\text{ and }
l\in\Hl{\varrho}{T},\text{ if }m(x)\in\Hl{\varrho}{T}.\end{equation}
Let \m{x\in N}, \m{k\in\HT{T}\cap\cgp{\cL}} and $y$ with \m{ky=x}, and let \m{l\in\Hl{\varrho}{T}} with \m{l^{-1}z=y};
we must show \m{\mH{\cL}(\cgp{\cL})\sum_\theta c_\theta\le\int\sum_{x'}G_{x'}(hl)\,d\mH{\cL}(h)}, where
\m{G_{x'}(l')=(\rePart\,v^*Rv)^{+}} for the vector \m{v_i=\corr{Z}{\omega_i}{l'}_{x'}} and \m{c_\theta=G_{n_\theta}(l_\theta)}.

The vector \m{l_ox=l_ok\,l^{-1}z} has vanishing coordinates at all $e_a$, $e^*_a$ and satisfies
\m{h(l_ox,l_ox)=h(z,z)}; so \m{m=m(x)\in\Hl{\varrho}{T}} and \m{\mathrm E(l_ox)^{-1}z=l_ox}, that is,
\m{m\,x=e}. For \m{\theta\in\Theta_0} the elements $l$ and \m{l_\theta mk} of \m{\Hl{\varrho}{T}} satisfy
$(l_\theta mk)^{-1}z=k^{-1}m^{-1}e=y=l^{-1}z$, so $b_\theta=l\,(l_\theta mk)^{-1}\in\HT{T}$ by Input~\ref{i:1}(d). Thus
\begin{equation}\label{e:l}l=b_\theta\,l_\theta\,m\,k\qquad(\theta\in\Theta_0).\end{equation}
Two invariances of the correlations: for \m{k'\in\HT{T}\cap\cgp{\cL}}, \m{l'\in\GT{T}} and a sequence $p$,
\m{\corr{Z}{p}{k'l'}=\corr{Z}{p}{l'}=\corr{Z}{p}{l'k'}}. Indeed \m{W(k'y')=W(y')} for all~$y'$, by
Lemma~\ref{l:phase}(d) with \m{u=1} and \m{\sd{1}=0}; this gives the first equality in the definition, and the
second after the substitution \m{y'\mapsto k'^{-1}y'}, since $k'$ commutes with \m{\weyl{\varrho}{p}} and
\m{\gM{\cL}} is left invariant.

Now let \m{h\in\HT{T}}, which carries \m{\mH{\cL}}, and let $\theta$ be such that \m{hb_\theta\in\cgp{\cL}}. By
\eqref{e:l}, the two invariances and \eqref{e:m},
\m{\corr{Z}{\omega_i}{hl}=\corr{Z}{\omega_i}{(hb_\theta)\,l_\theta\,m\,k}=\corr{Z}{\omega_i}{l_\theta}} for all~$i$, so
\m{G_{n_\theta}(hl)=c_\theta}. If also \m{hb_{\theta'}\in\cgp{\cL}}, then
\m{l_\theta l_{\theta'}^{-1}=b_\theta^{-1}b_{\theta'}\in\cgp{\cL}} by \eqref{e:l}; so \m{\theta\mapsto n_\theta} is injective
on the set of these~$\theta$, by the hypothesis. Hence
\[\sum_\theta\mathbf 1\bigl[hb_\theta\in\cgp{\cL}\bigr]\cdot c_\theta\ \le\ \sum_{x'}G_{x'}(hl).\]
Integrating over $h$ and using the invariance of \m{\mH{\cL}} under right translation by \m{b_\theta\in\HT{T}}, the
left side becomes \m{\mH{\cL}(\cgp{\cL})\sum_\theta c_\theta}.
\end{proof}

\sect{Where the results are used}{s:used}
All results are used inside the unit. Note~F (\emph{Translates of Poincaré series fill the rows}) quotes them
as its Input~5, for the form \m{T} in the lattice basis of the filling data over the product of the completions
at a finite set of places: Lemma~\ref{l:phase}(a), (c), (e) and Lemma~\ref{l:int}(c) (its part~(a)); Theorem~B and
Lemma~\ref{l:cong}~(b); Lemma~\ref{l:mass} and Lemma~\ref{l:int}(f)~(c); Theorem~A~(d); Theorem~C~(e);
Theorem~D(i)~(f); Theorem~D(ii)~(g); Lemma~\ref{l:cut}, Corollary~\ref{c:finite} and Lemma~\ref{l:avg}~(h). The
facts that this note recalls and does not prove, and that Note~F uses as well (the Haar property and the upper
bound for the mass of the integral group, the conditions on the measure on \m{\HT{T}}, the descent and the
bound of a floor form by its modulus), are Input~10 of Note~F, with the notes that own them. Note~G (\emph{From amplitude sums to local correlations})
uses Definition~\ref{d:phase}, Lemma~\ref{l:phase}, the finite Bessel vector and the correlations in its
unfolding theorem, Lemma~\ref{l:int}(a) and Lemma~\ref{l:int}(e), and the notion of an intermediate group.

\sect{Notes on the formal text}{s:notes}
\begin{enumerate}[leftmargin=1.6em,itemsep=3pt]
\item \emph{Arguments that the census reads differently from Lean; to be settled by} \texttt{\#check} \emph{on a
build.} The census reads the arguments from the text, and counts a variable of the file as an argument when
its name occurs in a declaration. For the energy (Definition~\ref{d:top}(c)) it lists five explicit arguments,
among them a sequence~$p$; in the definition $p$ is the index of the family, bound inside, and both uses of the
energy in the Lean text (the two statements of Theorem~D) apply it to four arguments: $Z$, $R$, $\omega$ and the
vector. The note writes these four. The same happens with the letter $r$ in Lemma~\ref{l:heads}, with $p$ in both
parts of Theorem~D, and with $u$ in Definition~\ref{d:phase}: each is bound in the statement and has no
counterpart as an argument. In the other direction, the census does not list $Z$, $p$, $q$ for Theorem~C(i),
nor $Z$ for Theorem~D(i), although the statements name these variables of the file; the note states them.
Implicit elements $x$, $y$, $k$, $g$ that the census lists for many statements are bound variables, or parts
of other names, and have no counterpart either.
\item \emph{Generality.} The form \m{\fzf{\cL}{Z}} is defined for every subgroup~$Z$ and every level; the
theorems take an intermediate group of a standard level. Lemma~\ref{l:phase} is stated in the formal text
without the hypothesis that the split Gram matrix is a hermitian unit.
\item \emph{The bound for the Bessel vector.} In Lemma~\ref{l:cut} the bound of the norm is \m{\sqrt{\#\Cos{Z}}},
from coordinates of modulus at most~$1$; the dossier's first reading said that the norm is at most~$1$, and
the second reader corrected it. Nothing depends on the value of the bound.
\item \emph{Arguments that a conclusion does not name.} In Corollary~\ref{c:finite} the number $r$ occurs only
in the hypothesis \m{r+10^8\nn{\fs}^4\le\kk{\fs}}, which thus says \m{10^8\nn{\fs}^4\le\kk{\fs}}. In Theorem~B the
number $t$ is free, so the bound holds with the largest $t$ allowed. In Proposition~\ref{p:depth} the
hypothesis is the existence of one point where the majorant does not vanish; it is used only to see that the
partial products of $p$ and $q$ are integral.
\item \emph{Hypotheses on the ring.} Theorem~D(ii) assumes that the units of $A$ form an open set and that
inversion is continuous at~$1$. Both hold for a finite product of local fields; they are hypotheses here
because the formal notion of a semilocal field is axiomatic. In Theorem~D(ii) the matrix $R$ is arbitrary.
\item \emph{Forms of the formal statements.} Proposition~\ref{p:shells} is stated in the formal text for a
gauge $\Gamma$ together with an equation that identifies it with a gauge of some depth; \m{\gix{\xi}^{t/2}} in
Theorem~A is written as a power of \m{\sqrt{\gix{\xi}}}, and the denominator index enters through its inverse;
the inequality \m{j\le\rh{\fs}} is needed to name the diagonal entry \m{\ZT{T}{x}_{jj}}. Differences of natural
numbers (the volume exponent, the exponents of the floor moduli) are truncated at~$0$ in the formal text. In
Theorem~D(i) the formal statement says \m{0\le\mathcal Q} for the order of the complex numbers, and bounds the
integral by the real part of~$\mathcal Q$.
\item \emph{Constants.} The exponents \m{10^8}, \m{10^9}, \m{10^{10}} are those of the formal text; the proofs
above show which smaller numbers are used: Proposition~\ref{p:shells} needs
\m{\vexp{\cL}\ge8\cdot10^7\nn{\fs}^4+t}, and Theorem~A applies it with \m{t+12}, which gives the sixth power of the
denominator index.
\item \emph{What was not re-derived.} The statements of Inputs 2 and~4 are the results of the tower method; this
note uses them as stated. The formal proof of Theorem~C(i) verifies the change of variables as an equality of
the laws of two functions on product spaces; the text above gives the substitutions. The measurability of the
functions that are integrated is checked in the formal text and not discussed here.
\end{enumerate}

\subhead{How this note was written.} From the structured proofs of the module on the top floor in the third
chunk of the unit's dossier, all of whose statements and claims a second reader had re-read; her remarks were
taken into account. I read the module in the Lean text line by line and compared every statement of this note
with it argument by argument; I read in the Lean text the definitions of the unit on the real place and of the
note on standard levels that are recalled in §\ref{s:setting}, and the statements of Inputs 2 to~4. Inputs 1
and~5 are stated from the dossier's readings and from the way the formal proofs use them; of these I read in the
Lean text the statements on primitive matrices, on Eichler elements, on the summation over ideals and the
definition of a valuation pair. Arguments that are mine and replace computations of the formal text: the proof
of Lemma~\ref{l:val} in terms of the lattice of the vectors $w$ with $Xw$ integral, the bookkeeping of the
exponents in Proposition~\ref{p:shells}, and the description of the change of variables in Theorem~C(i).
Nothing was compiled for this note. The program check of the unit's dossier covered \UnitSteps{} leaf steps of
\UnitProofs{} proofs.

\head{Appendix A. Concordance with the Lean artifact}
{\raggedright The formal names of the objects and the letters of this note: \texttt{phase} is \m{\phs{g}{z}};
\texttt{unitaryLevel}, \texttt{congGroup} are \m{\uLev{g}{T}}, \m{\cgp{\cL}}; \texttt{besselU}, \texttt{besselH},
\texttt{besselS}, \texttt{GT}, \texttt{unitriGroup} are \m{\UT{T}}, \m{\HT{T}}, \m{\ST{T}}, \m{\Hl{l}{T}},
\m{\Ql{l}{T}}; \texttt{simpleRootSum}, \texttt{frameGram}, \texttt{weylElem} are \m{\sd{u}}, \m{\ZT{T}{x}},
\m{\weyl{l}{p}}; \texttt{md}, \texttt{genIndex}, \texttt{denomIndex}, \texttt{idealZeta} are \m{\modA},
\m{\gix{u}}, \m{\den{\xi}}, \m{\zetaA{s}}; \texttt{groupMeasure}, \texttt{upperMeasure}, \texttt{floorMeasure},
\texttt{measureH}, \texttt{stackMeasure}, \texttt{floorQuadric} are \m{\gM{\cL}}, \m{\gMu{\cL}}, \m{\flm{\cL}{l}},
\m{\mH{\cL}}, \m{\stk{\cL}{l}{i}}, \m{\quadric{\cL}{\varrho}}; \texttt{ascConv heisenbergMeasure} is \m{\etab{j}};
\texttt{floorModulus}, \texttt{headDensity}, \texttt{volumeExponent}, \texttt{floorValues} are
\m{\fmod{\cL}{l}{g}}, \m{\hdens{\cL}{j}{p}{q}}, \m{\vexp{\cL}}, \m{\Val{\cL}{a}{x}{y}}; \texttt{upperPairing},
\texttt{pivotModulus}, \texttt{pairingNorm} are \m{\Av{\cL}{j}{k}{g}{y}}, \m{\Pv{\cL}{j}{k}{g}{y}}, \m{\PN{\cL}{W}{W'}};
\texttt{floorForm}, \texttt{formModulus}, \texttt{formMajorant}, \texttt{selfNormProduct} are
$\fform{\cL}{j}{p}{q}{W}{\omega}$ and its modulus, majorant and self-norm product; \texttt{besselCorrelation},
\texttt{finHeadWeight} are \m{\bcorr{\mu}{W}{\gamma}}, \m{\fhw{K}{j}{p}}; \texttt{IsStdLevel.gauge},
\texttt{Gauge.Λ}, \texttt{IsHead} are \m{\stdg{\cL}{\alpha}}, \m{\Lam{\Gamma}}, \m{\Hd{\Gamma}{j}{p}}; the level
\texttt{F} is~$\cL$, the index set \texttt{j} of Theorem~D is~$\mathcal C$, and \texttt{Level.energy} is
\m{\lenergy{Z}{R}{\omega}{e}}. Three statements as the artifact has them:\par}
\quoted{enormfloorFormle}
\quoted{floorZeroFormlemainterm}
\quoted{eventuallysumpacketsleenergy}
{\fontsize{8.6}{11.6}\selectfont\renewcommand{\leanfont}{\ttfamily\fontsize{7.7}{10}\selectfont\color{leancol}}
\setlength{\LTleft}{0pt}\setlength{\LTright}{\fill}\setlength{\tabcolsep}{0pt}
\begin{longtable}{@{}p{27mm}>{\raggedright\arraybackslash}p{68mm}p{34mm}r@{}}
\textbf{In this note} & \textbf{Lean declaration} & \textbf{File} & \textbf{Lines}\\\midrule\endhead
\ConcordanceRows
\end{longtable}}

\head{Notation}
\begingroup\fontsize{8.6}{10.9}\selectfont\raggedright
\setlength{\columnsep}{9mm}\setlength{\parskip}{0pt}
\begin{multicols}{2}
\nt{$A$, \m{x\mapsto\bar x}, $O_A$, \m{O_A^{\times}}, $\psi_A$, \m{\modA}, \m{\vol}}{a semilocal field with
involution, its integers, their units, the character, the modulus, Haar measure}
\nt{\m{\gix{u}}, \m{\den{\xi}}, \m{\zetaA{s}}, $I$, $I(x)$, $S(I)$}{generalised index, denominator index, ideal
zeta function; an ideal, the content ideal of a position, its shell}
\nt{\m{\fs=(\varrho,\kappa,k_0)}, \m{\rh{\fs}}, \m{\kk{\fs}}, \m{\nn{\fs}}, $n$}{a shape and its numbers; in proofs
\m{n=\nn{\fs}}, \m{\varrho=\rh{\fs}}}
\nt{\m{\Ix{\fs}}, $e_a$, $e^*_a$, $z$, \m{\mathfrak f_a}, $h$}{index set, basis vectors, frame, hermitian form}
\nt{$T$, \m{\JT{T}}, \m{\GT{T}}, \m{\ZT{T}{x}}}{an $\fs$-form, its split Gram matrix, unitary group, frame Gram
matrix}
\nt{\m{\UT{T}}, \m{\HT{T}}, \m{\ST{T}}, \m{\Hl{l}{T}}, \m{\Ql{l}{T}}}{Bessel unipotent subgroup, the subgroup
fixing the frame and its dual, Bessel subgroup, floor groups, unitriangular groups}
\nt{\m{\sd{u}}, \m{\weyl{l}{p}}, \m{\uLev{g}{T}}}{sum of the entries above the diagonal, Weyl element, congruence group
modulo~$g$}
\nt{$\iota$, $m$, $\cL$, \m{\Levels{\iota}{T}{m+1}}, \m{\gL{\cL}}, \m{\cgp{\cL}}}{index set of the orthonormal
part, number of further pairs, a level, the levels, its unit, its congruence subgroup}
\nt{\m{\BoxL{\cL}}, \m{\Posn{\cL}}, \m{\IsStd{\cL}}}{box, positions, standard levels}
\nt{$\eta_l$, \m{\etab{j}}, \m{\flm{\cL}{l}}, \m{\mH{\cL}}, \m{\stk{\cL}{l}{i}}, \m{\gM{\cL}}, \m{\gMu{\cL}},
\m{\quadric{\cL}{\varrho}}}{Heisenberg, floor, stack and group measures, the quadric measure; $\mu$, $\nu$, $\rho$,
$\theta$ measures}
\nt{\m{\fmod{\cL}{l}{g}}, \m{\hdens{\cL}{j}{p}{q}}, \m{\vexp{\cL}}, $s$, \m{\Val{\cL}{a}{x}{y}}}{floor modulus, head
density, volume exponent, floor values}
\nt{\m{\Av{\cL}{j}{k}{g}{y}}, \m{\Pv{\cL}{j}{k}{g}{y}}, \m{k^{d}_{W,W'}}, \m{\PN{\cL}{W}{W'}}}{averaged kernel,
pivot modulus, kernel of two functions, pairing norm}
\nt{\m{\fform{\cL}{j}{p}{q}{W}{\omega}}, \m{\Phi_j}, \m{\fmodu{\cL}{j}{p}{q}{W}{x}{\omega}},
\m{\fmaj{\cL}{j}{p}{q}{W}{W'}{x}}}{floor form, its integrand, modulus, majorant}
\nt{$\Gamma$, $\Gamma'$, $\Gamma''$, \m{\Gauges{\cL}}, \m{\szl{\Gamma}{x}{i}}, \m{\Lam{\Gamma}}, $\Lambda$,
\m{\Hd{\Gamma}{j}{p}}, \m{\stdg{\cL}{\alpha}}}{gauges, size, the constant, heads, gauge of depth~$\alpha$}
\nt{\m{\snp{\Gamma}{j}{p}{q}{W}{W'}{x}}, \m{\Sigma_0}, $V_0$}{self-norm product, its factor at floor~$0$, the
floor-mass constant}
\nt{\m{\fhw{K}{j}{p}}, \m{\bcorr{\mu}{W}{\gamma}}, \m{\ell^2(X)}}{head weight, Bessel correlation, the space \m{\CC^X}}
\nt{\m{\phs{g}{z}}, \m{\IntG{\cL}}, $Z$, \m{\Cos{Z}}, $a_c$, \m{\umass{\cL}}, $\mathsf v$, \m{\mr{Z}}}{finite
Bessel function, intermediate groups, cosets, representatives, unipotent mass, mass ratio}
\nt{\m{\fzf{\cL}{Z}}, \m{\bvec{Z}{y}}, $W$, \m{\pk{Z}}}{the form at the bottom, finite Bessel vector, its kernel}
\nt{\m{\corr{Z}{p}{l}}, \m{\topK{L}{Z}{p}{q}}, \m{\lenergy{Z}{R}{\omega}{e}}, $G$, $G_{x'}$}{correlation, top
kernel, energy, the integrands of the energy}
\nt{$p$, $q$, $r$, $P_l$, $Q_l$, $\omega_i$, $\mathcal C$, $R$, $\mathcal Q$}{sequences of units, partial
products, a family, its index set, a matrix, the kernel sum}
\nt{$\xi$, $\eta$, $\alpha$, $\beta$, $\gamma$, $\phi$, $\phi_0$, $\epsilon$, $\epsilon_0$,
$\varepsilon$}{units and real numbers of the proofs (frequency, its modification, depths, products, thresholds)}
\nt{$\omega$, \m{\varphi_\xi}, $F_\xi$, $\mathcal A$, $f$, $e$, $E$, $K$, $t$, $B$, $V$, $C$, $D$, $N$, $L$}{functions
and numbers of the statements; $C$ also a product of floor values or a compact set, $D$ a domain, $N$ a
neighbourhood, \m{L\subseteq A}}
\nt{$x$, $y$, $u$, $v$, $k$, $l$, $b$, $\sigma$, $\tau$, $w$, $d$, $c$, $X$, $Y_x$, $i$, $j$, $a$}{elements of
\m{\GT{T}}, vectors, matrices, units, indices, as said at each place}
\nt{$\Theta_0$, $\theta$, $l_\theta$, $n_\theta$, $v_\theta$, $c_\theta$, $b_\theta$, $o$, $m(x)$, $\mathrm E$,
$u_0$}{packets and the objects of the proof of Theorem~D}
\end{multicols}
\endgroup

\end{document}
